% concatenated from main.tex
\documentclass[12pt]{amsart}

%%% Preamble %%%
\usepackage{amsmath,amssymb,amsthm}
\usepackage{mathrsfs}
\usepackage{tikz-cd}
\usepackage{enumitem}
\usepackage{hyperref}
\usepackage{a4wide}
\setlength{\emergencystretch}{2em}

% Inequalities
\renewcommand{\leq}{\leqslant}
\renewcommand{\geq}{\geqslant}
\renewcommand{\le}{\leqslant}
\renewcommand{\ge}{\geqslant}

\theoremstyle{plain}
\newtheorem{theorem}{Theorem}[section]
\newtheorem{proposition}[theorem]{Proposition}
\newtheorem{lemma}[theorem]{Lemma}
\newtheorem{corollary}[theorem]{Corollary}

\theoremstyle{definition}
\newtheorem{definition}[theorem]{Definition}
\newtheorem{example}[theorem]{Example}

\theoremstyle{remark}
\newtheorem{remark}[theorem]{Remark}
\newtheorem{problem}[theorem]{Problem}
\newtheorem{notation}[theorem]{Notation}

\DeclareMathOperator{\FBL}{FBL}
\DeclareMathOperator{\SB}{SB}
\DeclareMathOperator{\Subl}{Subl}
\DeclareMathOperator{\Ideal}{Ideal}
\DeclareMathOperator{\dist}{dist}
\DeclareMathOperator{\Graph}{Graph}
\DeclareMathOperator{\WF}{WF}
\DeclareMathOperator{\Tr}{Tr}
\newcommand{\R}{\mathbb{R}}
\newcommand{\N}{\mathbb{N}}
\newcommand{\bbQ}{\mathbb{Q}}

\title{Polish spaces of separable Banach lattices}
\author[M.~Niwi\'nski]{Mariusz Niwi\'nski}
\address{Doctoral School of Exact and Natural Sciences, Faculty of Mathematics
and Computer Science, Institute of Mathematics, Jagiellonian University,
{\L}ojasiewicza 6, 30-348 Krak\'ow, Poland}
\email{mariusz.niwinski@doctoral.uj.edu.pl}
\date{May 2026}
\subjclass[2020]{46B42, 46E15, 03E15}
\keywords{Banach lattice, free Banach lattice, Polish space, Borel hierarchy, Wijsman topology, Michael's selection theorem}

\begin{document}

\begin{abstract}
We study the descriptive complexity of classes of separable Banach lattices.
Building on the theory of coding spaces for separable Banach spaces, we introduce
two Polish space encodings of separable Banach lattices: one via closed
sublattices of the universal lattice $\mathcal{C}=C(\Delta;L_1)$, and one via
closed order ideals of the free Banach lattice $\FBL[\ell_1]$.  We prove that,
for every separable Banach lattice $E$, the spaces of closed sublattices and of
closed order ideals of $E$ are Polish subspaces of the hyperspace of closed
subsets of $E$.  We also prove that the Fremlin projective tensor-product
operation on ideal codes is $\boldsymbol{\Sigma}^0_2$-measurable and has a
$G_\delta$ graph. % Finally, we discuss the relationship between the two coding
%spaces and record Borel upper bounds for several lattice-geometric properties,
%including uniform monotonicity and order uniform smoothness.
\end{abstract}

\maketitle

%=============================================================================
\section{Introduction}
%=============================================================================

This work studies the descriptive complexity of separable Banach lattices by
encoding them as elements of Polish spaces, extending classical results of
Bossard~\cite{Boss2002Cod} and Godefroy--Saint-Raymond~\cite{God2018Cod} from
the Banach space setting to the lattice setting.

The quotient-side coding used in this paper is the Banach-lattice instance of
Kania's quotient-encoding programme~\cite{Kan2026Pol}.  The conceptual source is
Kania's framework of separable quotient generators~\cite[Definition~3.11]{Kan2026Pol}:
objects are represented as quotients of a fixed generator, and definability is
read from the corresponding kernel data.  We use in particular the closedness and
Polishness theorem for admissible kernel spaces~\cite[Proposition~3.6]{Kan2026Pol},
the Wijsman-continuity of quotient-norm coordinates~\cite[Lemma~3.8]{Kan2026Pol},
and the quotient-code definability scheme~\cite[Main Theorem~3.13 and
Theorem~3.16]{Kan2026Pol}.  Kania also records that \(\FBL[\ell_1]\) is the
separable quotient generator for Banach lattices~\cite[Remark~5.1(c)]{Kan2026Pol}.
The sublattice coding via the universal lattice \(\mathcal C\) is complementary
to this quotient picture.

The paper is organised as follows.
In Section~\ref{sec:prelim-dst}, we recall the basic notions from descriptive
set theory and the theory of hyperspaces.
In Section~\ref{sec:prelim-lattice}, we collect the necessary background on
Banach lattices and free Banach lattices.
In Section~\ref{sec:coding}, we introduce two coding spaces of separable Banach
lattices and establish that the families of closed sublattices and closed order
ideals form Polish spaces.
In Section~\ref{sec:tensor}, we study the Borel complexity of the Fremlin tensor
product map.
%In Section~\ref{sec:equiv}, we discuss the relationship between the two coding
%spaces and formulate open problems.
In Section~\ref{sec:further}, we study several lattice-geometric properties
within the coding spaces: we establish $\boldsymbol{\Pi}^0_3$ bounds for uniform
monotonicity and order uniform smoothness (in the sense of
Kurc~\cite{Kurc1993Dual}), prove that strict monotonicity is coanalytic, restrict
the SICK-compactum discussion explicitly to the metrisable case, and formulate
the remaining completeness and sharpness problems.

%=============================================================================
\section{Preliminaries from descriptive set theory}\label{sec:prelim-dst}
%=============================================================================

We briefly recall the basic definitions.  A~\emph{Polish space} is a separable
and completely metrisable topological space.

\subsection{The Borel hierarchy}

Let $X$ be a Polish space.  The \emph{Borel hierarchy} consists of the families
$\boldsymbol{\Sigma}^0_\alpha$, $\boldsymbol{\Pi}^0_\alpha$, and
$\boldsymbol{\Delta}^0_\alpha$ for countable ordinals $\alpha \ge 1$,
defined inductively as follows.
\begin{enumerate}[label=(\roman*)]
    \item $\boldsymbol{\Sigma}^0_1$ is the family of open subsets of $X$.
    \item For each $\alpha$, the family $\boldsymbol{\Pi}^0_\alpha$ consists of
      the complements of sets in~$\boldsymbol{\Sigma}^0_\alpha$.
    \item For $\alpha > 1$, a set belongs to $\boldsymbol{\Sigma}^0_\alpha$ if
      and only if it can be written as a countable union of sets, each belonging
      to $\boldsymbol{\Pi}^0_\beta$ for some $\beta < \alpha$.
    \item $\boldsymbol{\Delta}^0_\alpha = \boldsymbol{\Sigma}^0_\alpha \cap
      \boldsymbol{\Pi}^0_\alpha$.
\end{enumerate}
In particular, $\boldsymbol{\Pi}^0_1$ is the family of closed sets,
$\boldsymbol{\Delta}^0_1$ is the family of clopen sets,
$\boldsymbol{\Sigma}^0_2$ is the family of $F_\sigma$~sets, and
$\boldsymbol{\Pi}^0_2$ is the family of $G_\delta$~sets.

A set $A \subseteq X$ is called
\emph{$\boldsymbol{\Sigma}^0_\alpha$-complete} if
$A \in \boldsymbol{\Sigma}^0_\alpha(X)$ and for every Polish space $Y$ and every
$B \in \boldsymbol{\Sigma}^0_\alpha(Y)$ there exists a continuous map
$f \colon Y \to X$ such that $f^{-1}(A) = B$.  The notion of
\emph{$\boldsymbol{\Pi}^0_\alpha$-complete} is defined analogously.
In particular, a $\boldsymbol{\Sigma}^0_\alpha$-complete set does not belong to
$\boldsymbol{\Pi}^0_\alpha$, hence not to $\boldsymbol{\Delta}^0_\alpha$.

\subsection{Hyperspaces of closed sets}\label{subsec:hyper}

Let $X$ be a compact metric space.  Denote by $\mathcal{F}(X)$ the family of
all non-empty closed subsets of $X$, equipped with the \emph{Vietoris topology},
which is generated by the subbase consisting of
\[
    E^+(U) = \{F \in \mathcal{F}(X) \colon F \cap U \neq \emptyset\}, \qquad
    E^-(U) = \{F \in \mathcal{F}(X) \colon F \subseteq U\},
\]
for all open $U \subseteq X$.  This topology coincides with the one induced by
the \emph{Hausdorff--Pompeiu metric}
\[
    d_H(E,F) = \max\Bigl\{\sup_{x \in E}\dist(x,F),\;\sup_{y \in F}\dist(y,E)\Bigr\}.
\]
When $X$ is compact, $\bigl(\mathcal{F}(X), d_H\bigr)$ is a Polish space.

If $X$ is merely Polish (but not compact), $\mathcal{F}(X)$ with the Vietoris
topology is no longer separable, and there is no canonical Polish topology on it.
Following~\cite{God2018Cod}, we call a topology $\tau$ on $\mathcal{F}(X)$
\emph{admissible} if $\bigl(\mathcal{F}(X), \tau\bigr)$ is a Polish space and:
\begin{enumerate}[label=(A\arabic*)]
    \item\label{A1} For every open $U \subseteq X$, the set $E^+(U)$ is
      $\tau$-open.
    \item\label{A2} There exists a $\tau$-subbase each element of which is a
      countable union of sets of the form $E^+(U) \setminus E^+(V)$ for open
      $U, V \subseteq X$.
\end{enumerate}
The Borel $\sigma$-algebra generated by any admissible topology equals the
$\sigma$-algebra generated by $\{E^+(U) \colon U \subseteq X \text{ open}\}$,
known as the \emph{Effros--Borel $\sigma$-algebra};
see~\cite[Section~2]{God2018Cod}.
Moreover, if $\tau_1$ and $\tau_2$ are two admissible topologies, then every
$\tau_1$-open set is $\boldsymbol{\Sigma}^0_2$ in~$\tau_2$.

It is natural to impose a further compatibility condition:
\begin{enumerate}[resume*]
    \item\label{A3} The set
      $\{(x, F) \in X \times \mathcal{F}(X) \colon x \in F\}$ is closed in the
      product topology on $X \times \mathcal{F}(X)$.
\end{enumerate}
Note that the membership relation $x \in F$ is always
$\boldsymbol{\Pi}^0_2$ with respect to any admissible topology: fixing a
countable base $\{U_n\}_{n \in \omega}$ of $X$, one has
\[
    \{(x,F) \colon x \in F\}
    = \bigcap_{n \in \omega}
      \Bigl[(X \setminus U_n) \times \mathcal{F}(X) \;\cup\; X \times E^+(U_n)\Bigr].
\]

\begin{example}\label{ex:compact}
Since every Polish space $X$ is homeomorphic to a totally bounded metric space,
we may choose a compatible metric $\rho$ under which $X$ is totally bounded.  Let
$\hat{X}$ denote the completion of $(X, \rho)$; then $\hat{X}$ is compact.
Identify $\mathcal{F}(X)$ with
\[
    \mathcal{F}^*(X) = \bigl\{K \in \mathcal{F}(\hat{X}) \colon
    \overline{X \cap K}^{\hat{X}} = K\bigr\},
\]
which is a $\boldsymbol{\Pi}^0_2$~subset of the Polish space
$\mathcal{F}(\hat{X})$.  By Alexandrov's theorem, $\mathcal{F}^*(X)$ is itself
Polish.  The restriction of the Vietoris topology of $\mathcal{F}(\hat{X})$ to
$\mathcal{F}^*(X)$ is admissible and satisfies~\ref{A3};
see, e.g.,~\cite[Chapter~2]{Dod2010Ban}.
\end{example}

\begin{example}[Wijsman topology]\label{ex:wijsman}
Let $d$ be a compatible metric on a Polish space $X$, and fix a dense sequence
$\{\alpha_n\}_{n \in \omega} \subseteq X$.  The map
\[
    \mathcal{F}(X) \ni F \;\longmapsto\; \bigl(d(\alpha_n, F)\bigr)_{n \in \omega}
    \in \R^\omega
\]
is injective.  The \emph{Wijsman topology} is the coarsest topology making this
map a homeomorphism onto its image, where $\R^\omega$ carries the product
topology.  It is shown in~\cite{God2018Cod} that the Wijsman topology is
admissible and independent of the choice of dense sequence.  Moreover, the
evaluation map
\[
    X \times \mathcal{F}(X) \ni (x, F) \longmapsto d(x, F)
\]
is continuous when $\mathcal{F}(X)$ carries the Wijsman topology.  Indeed,
suppose $F_n \to F$ in the Wijsman topology and fix $x \in X$.  Given
$\varepsilon > 0$, choose $m$ with $d(x, \alpha_m) < \varepsilon$.  Using the
$1$-Lipschitz estimate $\lvert d(x, G) - d(\alpha_m, G) \rvert \le d(x, \alpha_m)$
for $G \in \mathcal{F}(X)$, we obtain
\[
    \lvert d(x, F_n) - d(x, F) \rvert
    \le 2\varepsilon + \lvert d(\alpha_m, F_n) - d(\alpha_m, F) \rvert,
\]
so $d(x, F_n) \to d(x, F)$ as $\varepsilon \downarrow 0$.  Joint continuity in
$(x, F)$ then follows from $\lvert d(x, F) - d(x', F) \rvert \le d(x, x')$.
Consequently,
\[
    \{(x, F) \in X \times \mathcal{F}(X) \colon x \in F\}
    = \{(x, F) \colon d(x, F) = 0\}
\]
is closed, so condition~\ref{A3} holds.  For closed subspaces or ideals, the
same distance coordinate is the quotient norm \(F\mapsto\|x+F\|\); this is the
basic Wijsman-continuity mechanism used in Kania's framework
\cite[Lemma~3.8]{Kan2026Pol}.
\end{example}

\begin{lemma}\label{lem:closure-map-borel}
Let $X$ be a Polish space and endow $\mathcal{F}(X)$ with its Effros--Borel
$\sigma$-algebra (equivalently, with the Borel $\sigma$-algebra of any admissible
topology).  Define
\[
    \mathrm{Cl}_X \colon X^\omega \to \mathcal{F}(X), \qquad
    \mathrm{Cl}_X\bigl((x_n)_{n \in \omega}\bigr)
    = \overline{\{x_n \colon n \in \omega\}}.
\]
Then $\mathrm{Cl}_X$ is Borel measurable.
\end{lemma}

\begin{proof}
Let $U \subseteq X$ be open.  Since $U$ is open,
$\overline{\{x_n \colon n \in \omega\}} \cap U \ne \emptyset$ if and only if
$x_n \in U$ for some~$n$.  Hence
\[
    \mathrm{Cl}_X^{-1}\bigl(E^+(U)\bigr)
    = \bigcup_{n \in \omega}
      \bigl\{(x_j)_{j \in \omega} \in X^\omega \colon x_n \in U\bigr\},
\]
which is open in $X^\omega$.  The sets $E^+(U)$ generate the Effros--Borel
$\sigma$-algebra, so $\mathrm{Cl}_X$ is Borel.
\end{proof}

%=============================================================================
\section{Preliminaries on Banach lattices}\label{sec:prelim-lattice}
%=============================================================================

Throughout, all Banach spaces are over~$\R$.

\subsection{Basic definitions}

\begin{definition}
A \emph{Banach lattice} is a Banach space $(E, \lVert\cdot\rVert_E)$ equipped
with a lattice order~$\preceq$ (that is, a partial order in which every pair
$x, y \in E$ has a supremum $x \vee y$ and an infimum $x \wedge y$) satisfying,
for all $x, y, z \in E$ and $\alpha > 0$:
\begin{enumerate}[label=(\roman*)]
    \item $x \preceq y \implies x + z \preceq y + z$,
    \item $x \preceq y \implies \alpha x \preceq \alpha y$,
    \item $\lvert x \rvert \preceq \lvert y \rvert \implies
      \lVert x \rVert_E \le \lVert y \rVert_E$,
\end{enumerate}
where $\lvert x \rvert = x \vee (-x)$ denotes the \emph{modulus} of~$x$.
\end{definition}

A \emph{Banach lattice homomorphism} is a linear operator between Banach lattices
that preserves the lattice operations (equivalently, $T(\lvert x \rvert)
= \lvert Tx \rvert$ for all~$x$).  Every such homomorphism is automatically
continuous.

\begin{lemma}\label{lem:lattice-ops-lipschitz}
Let $E$ be a Banach lattice.  Then:
\begin{enumerate}[label=\textup{(\roman*)}]
    \item $\bigl\lVert\,\lvert x \rvert - \lvert y \rvert\,\bigr\rVert_E
      \le \lVert x - y \rVert_E$ for all $x, y \in E$;
    \item the lattice operations $\vee$ and $\wedge$ are Lipschitz:
      \[
          \lVert (x \vee y) - (x' \vee y') \rVert_E
          \le \lVert x - x' \rVert_E + \lVert y - y' \rVert_E,
      \]
      and similarly for $\wedge$.
\end{enumerate}
In particular, $x \mapsto \lvert x \rvert$ and
$(x,y) \mapsto x \vee y,\; x \wedge y$ are continuous.
\end{lemma}

\begin{proof}
The lattice inequality
$\bigl\lvert\, \lvert x \rvert - \lvert y \rvert \,\bigr\rvert
\preceq \lvert x - y \rvert$ holds in every vector lattice.
Applying the norm monotonicity axiom yields~(i).  For~(ii), use the identities
\[
    x \vee y = \tfrac{1}{2}\bigl(x + y + \lvert x - y \rvert\bigr), \qquad
    x \wedge y = \tfrac{1}{2}\bigl(x + y - \lvert x - y \rvert\bigr),
\]
and combine~(i) with the triangle inequality.
\end{proof}

\begin{example}\label{ex:classical-lattices}
The following are standard examples of Banach lattices, each with the pointwise
(respectively, pointwise a.e.)\ order:
\begin{enumerate}[label=(\alph*)]
    \item $C(K; \R)$ for any compact Hausdorff space~$K$;
    \item $L_p(X, \mathcal{A}, \mu; \R)$ for $p \in [1, \infty)$ and a
      $\sigma$-finite measure space $(X, \mathcal{A}, \mu)$;
    \item $\ell_p(\omega; \R)$ for $p \in [1, \infty)$.
\end{enumerate}
\end{example}

\begin{definition}\label{def:sublattice-ideal}
Let $E$ be a Banach lattice.
\begin{enumerate}[label=(\roman*)]
    \item A closed linear subspace $F \subseteq E$ is a \emph{sublattice} if
      $\lvert y \rvert \in F$ for every $y \in F$.
    \item A sublattice $J \subseteq E$ is an \emph{order ideal} (or simply
      \emph{ideal}) if it is \emph{solid}: whenever $x, y \in E$ satisfy
      $\lvert x \rvert \preceq \lvert y \rvert$ and $y \in J$, one has $x \in J$.
    \item For a non-zero element $u \succeq \mathbf{0}_E$, the \emph{principal
      ideal} generated by~$u$ is
      $J_u = \{y \in E \colon \exists\, r > 0 \text{ with }
      \lvert y \rvert \preceq r u\}$.
    \item The \emph{positive cone} is
      $E^+ = \{x \in E \colon x \succeq \mathbf{0}_E\}$.
    \item $E$ is \emph{Dedekind complete} if every non-empty, order-bounded
      subset of $E$ has a supremum in~$E$.  When we refer to the
      \emph{Dedekind completion} $\hat E$ of a Banach lattice, we mean the
      usual Dedekind completion of the underlying Archimedean vector lattice,
      equipped with the lattice norm used in the cited source.  In this paper
      $\hat E$ is used only through quoted results on the Fatou and Nakano
      properties.
\end{enumerate}
\end{definition}

We work with the following two categories:
\begin{itemize}
    \item $\mathfrak{ban}_s$: separable real Banach spaces with contractive
      (i.e., norm~$\le 1$) linear operators as morphisms.
    \item $\mathfrak{banlat}_s$: separable real Banach lattices with contractive
      lattice homomorphisms as morphisms.
\end{itemize}

\subsection{Free Banach lattices}\label{subsec:FBL}

For every Banach space $X$, the \emph{free Banach lattice generated by~$X$},
denoted $\FBL[X]$, is characterised by the following universal property: there
exists an isometric embedding $\delta_X \colon X \to \FBL[X]$ such that, for
every Banach lattice $E$ and every contractive linear operator
$T \colon X \to E$, there exists a unique contractive lattice homomorphism
$\hat{T} \colon \FBL[X] \to E$ satisfying $\hat{T} \circ \delta_X = T$:
\[
    \begin{tikzcd}
    \FBL[X] \arrow[rd, "\hat{T}"]  & \\
    X \arrow[r, "T"'] \arrow[u, "\delta_X"] & E
    \end{tikzcd}
\]
Moreover, $\lVert \hat{T} \rVert = \lVert T \rVert$.

The assignment $X \mapsto \FBL[X]$ extends to a covariant functor from
$\mathfrak{ban}_s$ to $\mathfrak{banlat}_s$.

An explicit construction is given by Avil\'es, Rodr\'\i guez, and
Tradacete~\cite{aviles2018free}.  Let
\[
    H[X] = \{f \colon X' \to \R \colon f \text{ is positively homogeneous}\},
\]
equipped with the pointwise order, and define
\[
    \lVert f \rVert_{\FBL[X]}
    = \sup\biggl\{\sum_{i=1}^m \lvert f(x_i^*) \rvert \colon
      m \in \N,\; \{x_i^*\}_{i=1}^m \subseteq X',\;
      \sup_{x \in \overline{B}_X} \sum_{i=1}^m \lvert x_i^*(x) \rvert
      \le 1 \biggr\},
\]
where $\overline{B}_X$ is the closed unit ball of~$X$.  Set
$H_1[X] = \{f \in H[X] \colon \lVert f \rVert_{\FBL[X]} < \infty\}$.
The evaluation map $\delta_X \colon X \to H_1[X]$, defined by
$\delta_X(x)(x^*) = x^*(x)$, is a linear isometry, and $\FBL[X]$ is the closure
of the sublattice generated by $\delta_X(X)$ in~$H_1[X]$.

The following facts are proved in~\cite{aviles2018free}:
\begin{enumerate}[label=(\alph*)]
    \item $\FBL[\ell_1(\omega; \R)]$ is isometrically lattice-isomorphic to
      $\FBL[\omega]$, the free Banach lattice generated by the set~$\omega$.
    \item $\FBL[\ell_1(\omega; \R)]$ is generated (as a Banach lattice) by the
      countable set
      $\{\delta_{e_j}\}_{j \in \omega}$, where $\{e_j\}_{j \in \omega}$ is the
      standard basis of~$\ell_1$.
    \item Every separable Banach lattice $E$ is isometrically lattice-isomorphic
      to a quotient $\FBL[\ell_1(\omega; \R)]/J$ for some closed ideal~$J$.
      More explicitly, choose a norm-one quotient map
      $q\colon \ell_1(\omega;\R)\to E$ such that, for every $y\in E$ and
      every $\eta>0$, there is $x\in\ell_1$ with $q x=y$ and
      $\|x\|\le \|y\|+\eta$.  The universal property gives a contractive
      lattice homomorphism $\widehat q\colon \FBL[\ell_1]\to E$ extending
      $q$.  Since $q(\ell_1)=E$, the map $\widehat q$ is onto; if
      $J=\ker \widehat q$, then the quotient norm on
      $\FBL[\ell_1]/J$ agrees with the given norm of $E$.
\end{enumerate}
Consequently, \(\FBL[\ell_1]\) is the separable quotient generator for the
Banach-lattice category in the sense of Kania's Definition~3.11; compare
\cite[Remark~5.1(c)]{Kan2026Pol}.

\subsection{Universal embedding}

Let $\Delta = \{0,1\}^\omega$ denote the Cantor set.
By a result of Leung, Li, Oikhberg, and Tursi~\cite{LLOT2019Uni}, the Banach
lattice
\[
    C\bigl(\Delta;\, L_1(0,1; \R)\bigr),
\]
with the supremum norm and pointwise a.e.\ order, is injectively universal for
separable Banach lattices: every separable Banach lattice admits a lattice
isometric embedding as a closed sublattice of this space.

\begin{notation}
For brevity, we write $\mathcal{C} = C(\Delta; L_1)$, where
$L_1 = L_1(0,1; \R)$, and $\FBL[\ell_1]$ for $\FBL[\ell_1(\omega; \R)]$.
\end{notation}

%=============================================================================
\section{Coding spaces of separable Banach lattices}\label{sec:coding}
%=============================================================================

\subsection{Spaces of closed subspaces}\label{subsec:SB}

Let $X$ be a separable Banach space.  Denote by $\SB(X) \subseteq \mathcal{F}(X)$
the set of all closed linear subspaces of~$X$, where $\mathcal{F}(X)$ carries an
admissible topology satisfying~\ref{A1}--\ref{A3}
(for instance, the Wijsman topology).

The following is proved in~\cite{God2018Cod} (see
also~\cite[Chapter~2]{Dod2010Ban}).

\begin{proposition}\label{prop:SB-Polish}
$\SB(X)$ is a $\boldsymbol{\Pi}^0_2$~subset of $\mathcal{F}(X)$, and hence a
Polish space in the subspace topology.
\end{proposition}

A key tool is the following consequence of Michael's continuous selection theorem;
see~\cite[Theorem~4.1]{God2018Cod}.

\begin{theorem}[Godefroy--Saint-Raymond]\label{thm:selection}
Let $X$ be a separable Banach space.  There exist continuous functions
$f_j \colon \SB(X) \to X$, $j \in \omega$, such that for every
$F \in \SB(X)$ one has
\[
    f_j(F) \in F \qquad (j \in \omega),
\]
and
\[
    F = \overline{\{f_j(F) \colon j \in \omega\}}.
\]
\end{theorem}

\subsection{Closed sublattices}\label{subsec:subl}

Let $E$ be a separable Banach lattice and let $f_j \colon \SB(E) \to E$,
$j \in \omega$, be the continuous selections from Theorem~\ref{thm:selection}.

\begin{proposition}\label{prop:subl}
A closed subspace $F \in \SB(E)$ is a sublattice of~$E$ if and only if
\begin{equation}\label{eq:subl-cond}
    \lvert f_j(F) \rvert \in F \qquad \text{for all } j \in \omega.
\end{equation}
Moreover, $\Subl(E)$ is a closed ($\boldsymbol{\Pi}^0_1$) subset of $\SB(E)$
and hence a Polish space.
\end{proposition}

\begin{proof}
If $F$ is a sublattice then~\eqref{eq:subl-cond} holds trivially.  Conversely,
suppose~\eqref{eq:subl-cond} holds.  For any $y \in F$, choose indices
$(j_n)_{n \in \omega}$ such that $f_{j_n}(F) \to y$.  By
Lemma~\ref{lem:lattice-ops-lipschitz}\,(i) and~\eqref{eq:subl-cond}, the
corresponding moduli lie in~$F$ and converge to $\lvert y \rvert$, so
$\lvert y \rvert \in F$ by closedness.

For the descriptive complexity, write
\[
    \Subl(E) = \bigcap_{j \in \omega}
    \bigl\{F \in \SB(E) \colon \lvert f_j(F) \rvert \in F\bigr\}.
\]
The map $F \mapsto \bigl(\lvert f_j(F) \rvert,\, F\bigr) \in E \times \SB(E)$
is continuous (using Lemma~\ref{lem:lattice-ops-lipschitz}\,(i) and continuity
of~$f_j$), and condition~\ref{A3} ensures that
$\{(x, F) \colon x \in F\}$ is closed.  Hence each set in the intersection is
closed, and their intersection is closed.
\end{proof}

Taking $E = \mathcal{C}$ yields a Polish space
$\Subl(\mathcal{C})$ that serves as a coding space for
separable Banach lattices: every separable Banach lattice is isometrically
lattice-isomorphic to some element of this space.

\subsection{Closed order ideals}\label{subsec:ideal}

Kania explicitly notes that Banach lattices fit the quotient-kernel framework,
and his general admissible-kernel theorem applies after taking \(\vee\),
\(\wedge\), and \(|\cdot|\) as the continuous operations and closed lattice ideals
as kernels; see \cite[Remark~3.5 and Proposition~3.6]{Kan2026Pol}.  The next
theorem gives a direct lattice proof in our notation and, in particular, records
the explicit solidity test used later.

\begin{theorem}\label{thm:ideal}
Let $E$ be a separable Banach lattice and let
$f_j \colon \SB(E) \to E$, $j \in \omega$, be the continuous selections from
Theorem~\ref{thm:selection}.  Fix a dense sequence $(g_k)_{k \in \omega}$
in~$E$.

Then $F \in \Ideal(E)$ if and only if $F \in \SB(E)$ and:
\begin{enumerate}[label=\textup{(I\arabic*)}]
    \item\label{I1} $\lvert f_j(F) \rvert \in F$ for all $j \in \omega$
      \quad\textup{(}sublattice condition\textup{)};
    \item\label{I2} $\lvert g_k \rvert \wedge \lvert f_j(F) \rvert \in F$
      for all $j, k \in \omega$
      \quad\textup{(}solidity tested on dense elements\textup{)}.
\end{enumerate}
Moreover, $\Ideal(E)$ is a closed ($\boldsymbol{\Pi}^0_1$) subset of $\SB(E)$
and hence a Polish space.
\end{theorem}

\begin{proof}
\emph{Necessity.}
If $F$ is a closed ideal then it is, in particular, a closed sublattice,
so~\ref{I1} holds.
For~\ref{I2}, note that
$\lvert g_k \rvert \wedge \lvert f_j(F) \rvert \preceq \lvert f_j(F) \rvert$
and $f_j(F) \in F$, hence solidity gives
$\lvert g_k \rvert \wedge \lvert f_j(F) \rvert \in F$.

\medskip
\emph{Sufficiency.}
Assume~\ref{I1} and~\ref{I2}.  By Proposition~\ref{prop:subl},
condition~\ref{I1} implies that $F$ is a closed sublattice of~$E$.

We verify solidity.  Let $x \in E$ and $y \in F$ satisfy
$\lvert x \rvert \preceq \lvert y \rvert$.  Set $v := \lvert y \rvert \in F^+$
(the modulus lies in $F$ by~\ref{I1}).  It suffices to show that every
$u \in E^+$ with $u \preceq v$ belongs to~$F$, because then we may apply this
to $u = x^+$ and $u = x^-$ (both satisfy
$\mathbf{0}_E \preceq x^\pm \preceq \lvert x \rvert \preceq v$) and conclude
$x = x^+ - x^- \in F$.

Let $u \in E^+$ with $u \preceq v$.  Choose indices $(k_n)_{n \in \omega}$ such that
$g_{k_n} \to u$ in norm.  Then $\lvert g_{k_n} \rvert \to \lvert u \rvert = u$
by Lemma~\ref{lem:lattice-ops-lipschitz}\,(i).  Choose also indices
$(j_m)_{m \in \omega}$ such that
$f_{j_m}(F) \to v$; then
$\lvert f_{j_m}(F) \rvert \to \lvert v \rvert = v$.

By~\ref{I2}, we have $\lvert g_{k_n} \rvert \wedge \lvert f_{j_m}(F) \rvert \in F$
for all $m, n$.  Fixing $n$ and letting $m \to \infty$, the Lipschitz continuity
of $\wedge$ (Lemma~\ref{lem:lattice-ops-lipschitz}\,(ii)) and closedness
of~$F$ yield $\lvert g_{k_n} \rvert \wedge v \in F$.  Letting $n \to \infty$,
we have $\lvert g_{k_n} \rvert \wedge v \to u \wedge v$ by continuity of $\wedge$
(Lemma~\ref{lem:lattice-ops-lipschitz}\,(ii)).  Since $F$ is closed and each
$\lvert g_{k_n} \rvert \wedge v$ belongs to~$F$, it follows that
$u \wedge v \in F$.  Since $u \preceq v$, we have
$u \wedge v = u$, hence $u \in F$.

\medskip
\emph{Descriptive complexity.}
We may write
\[
    \Ideal(E)
    = \bigcap_{j \in \omega}
      \bigl\{F \in \SB(E) \colon \lvert f_j(F) \rvert \in F\bigr\}
    \;\cap\;
    \bigcap_{j,k \in \omega}
      \bigl\{F \in \SB(E) \colon
      \lvert g_k \rvert \wedge \lvert f_j(F) \rvert \in F\bigr\}.
\]
Each set in both families is closed in~$\SB(E)$: the map
$F \mapsto \bigl(\lvert f_j(F) \rvert, F\bigr)$ (respectively,
$F \mapsto \bigl(\lvert g_k \rvert \wedge \lvert f_j(F) \rvert, F\bigr)$) is
continuous by Lemma~\ref{lem:lattice-ops-lipschitz} and continuity of~$f_j$,
and~\ref{A3} makes the membership relation $\{(x, F) \colon x \in F\}$ closed.
Hence $\Ideal(E)$ is a countable intersection of closed sets, which is closed.
\end{proof}

Theorem~\ref{thm:ideal} is the Banach-lattice version of the admissible-kernel
closedness theorem in Kania's quotient framework~\cite[Proposition~3.6]{Kan2026Pol}.
The proof above is included because it identifies the abstract kernel condition
with the concrete closed-order-ideal tests used later.

\begin{remark}\label{rem:positive-dense}
Since $E^+$ is a closed subset of a separable metric space, it is separable.
Thus one may, if desired, fix the dense sequence $(g_k)$ in
Theorem~\ref{thm:ideal} inside~$E^+$.  In that case condition~\ref{I2} reads
simply $g_k \wedge \lvert f_j(F) \rvert \in F$ for all $j, k$.
\end{remark}

\begin{corollary}\label{cor:coding-Polish}
For any separable Banach lattice $E$, the space $\Subl(E)$ is a closed
($\boldsymbol{\Pi}^0_1$) subset of $\SB(E)$, and $\Ideal(E)$ is a closed
($\boldsymbol{\Pi}^0_1$) subset of $\SB(E)$.  In particular, both are
Polish spaces.
\end{corollary}

\begin{corollary}\label{cor:ambient-complexity}
Let $E$ be a separable Banach lattice and endow $\mathcal{F}(E)$ with any
admissible topology.  Then $\Subl(E)$ and $\Ideal(E)$ are
$\boldsymbol{\Pi}^0_2$ subsets of $\mathcal{F}(E)$.
\end{corollary}

\begin{proof}
Let $\tau_W$ denote the Wijsman topology on $\mathcal{F}(E)$, and let $\tau$
be an arbitrary admissible topology on $\mathcal{F}(E)$.

By Proposition~\ref{prop:subl} and Theorem~\ref{thm:ideal}, the sets
$\Subl(E)$ and $\Ideal(E)$ are closed in $\SB(E)$ when $\SB(E)$ is viewed as a
subspace of $(\mathcal{F}(E),\tau_W)$.  Hence there exist $\tau_W$-closed sets
$C_{\Subl}, C_{\Ideal} \subseteq \mathcal{F}(E)$ such that
\[
    \Subl(E)=\SB(E)\cap C_{\Subl},
    \qquad
    \Ideal(E)=\SB(E)\cap C_{\Ideal}.
\]
For instance, one may take
$C_{\Subl}=\overline{\Subl(E)}^{\,\tau_W}$ and
$C_{\Ideal}=\overline{\Ideal(E)}^{\,\tau_W}$.

By the comparison theorem for admissible topologies
(\cite[Section~2]{God2018Cod}), every $\tau_W$-open set is
$\boldsymbol{\Sigma}^0_2$ with respect to $\tau$.  Equivalently, the identity
map
\[
    \mathrm{id}\colon (\mathcal{F}(E),\tau)\longrightarrow
    (\mathcal{F}(E),\tau_W)
\]
is of Baire class~$1$.  Therefore the preimage of every $\tau_W$-closed set is
a $\boldsymbol{\Pi}^0_2$ subset of $(\mathcal{F}(E),\tau)$.

In particular, both $C_{\Subl}$ and $C_{\Ideal}$ are
$\boldsymbol{\Pi}^0_2$ in $(\mathcal{F}(E),\tau)$.  Since $\SB(E)$ is
$\boldsymbol{\Pi}^0_2$ in $(\mathcal{F}(E),\tau)$ by
Proposition~\ref{prop:SB-Polish}, it follows that
\[
    \Subl(E)=\SB(E)\cap C_{\Subl},
    \qquad
    \Ideal(E)=\SB(E)\cap C_{\Ideal}
\]
are $\boldsymbol{\Pi}^0_2$ subsets of $(\mathcal{F}(E),\tau)$.
\end{proof}

Taking $E = \FBL[\ell_1]$, the space $\Ideal\bigl(\FBL[\ell_1]\bigr)$ is a
Polish space that encodes all separable Banach lattices (up to isometric lattice
isomorphism) via the quotient representation $E \cong \FBL[\ell_1] / J$.
This is precisely the Banach-lattice quotient coding of
\cite[Main Theorem~3.13 and Remark~5.1(c)]{Kan2026Pol}: the Wijsman coordinates
$J\mapsto\|x+J\|$ are continuous, and the full quotient-norm code is faithful as
a code of kernels, although different kernels may still give isomorphic
quotients.  Similarly, taking $E = \mathcal{C}$, the space
$\Subl(\mathcal{C})$ provides an alternative embedding-based coding.

\subsection{A Borel ideal-hull operator}\label{subsec:ideal-hull}

Let $E$ be a separable Banach lattice, fix continuous selections
$f_j \colon \SB(E) \to E$ as in Theorem~\ref{thm:selection}, and fix a dense
sequence $(g_k)_{k \in \omega}$ in~$E$.

\begin{proposition}\label{prop:ideal-hull}
For $F \in \Subl(E)$ define
\[
    \Gamma(F)
    := \overline{\mathrm{span}}\Bigl\{\,
    \lvert g_k \rvert \wedge \lvert f_j(F) \rvert
    \colon j, k \in \omega \,\Bigr\}.
\]
Then $\Gamma(F)$ is the smallest closed ideal of $E$ containing~$F$
(i.e.\ the closed ideal generated by~$F$).  Moreover, the map
$\Gamma \colon \Subl(E) \to \Ideal(E)$ is Borel measurable.
\end{proposition}

\begin{proof}
Let $\langle F \rangle_{\Ideal}$ denote the closed ideal generated by~$F$,
i.e.\ the norm-closure of the (not necessarily closed) ideal
\[
    I(F) := \{x \in E \colon \exists\, y \in F \text{ with }
    \lvert x \rvert \preceq \lvert y \rvert\}.
\]

\emph{$\Gamma(F) \subseteq \langle F \rangle_{\Ideal}$.}
For each $j$ we have $f_j(F) \in F$, hence $\lvert f_j(F) \rvert \in F$ because
$F$ is a sublattice.  Since $\langle F \rangle_{\Ideal}$ is an ideal containing
$F$ and $\lvert g_k \rvert \wedge \lvert f_j(F) \rvert \preceq \lvert f_j(F) \rvert$,
solidity gives
$\lvert g_k \rvert \wedge \lvert f_j(F) \rvert \in \langle F \rangle_{\Ideal}$
for all $j, k$.  Thus $\Gamma(F) \subseteq \langle F \rangle_{\Ideal}$.

\medskip
\emph{$I(F) \subseteq \Gamma(F)$.}
Let $x \in I(F)$, so $\lvert x \rvert \preceq v$ for some $v \in F^+$.
Write $x = x^+ - x^-$, where $x^\pm \in E^+$ and
$\mathbf{0}_E \preceq x^\pm \preceq \lvert x \rvert \preceq v$.  Since
$\Gamma(F)$ is a linear subspace, it suffices to show that every $u \in E^+$
with $u \preceq v$ belongs to $\Gamma(F)$.

Fix such~$u$.  Choose indices $(k_n)_{n \in \omega}$ such that $g_{k_n} \to u$ in norm.
Then $\lvert g_{k_n} \rvert \to u$ by
Lemma~\ref{lem:lattice-ops-lipschitz}\,(i).  Also choose indices
$(j_m)_{m \in \omega}$ such that $f_{j_m}(F) \to v$; then
$\lvert f_{j_m}(F) \rvert \to v$.

For each fixed $n$ we have
$\lvert g_{k_n} \rvert \wedge \lvert f_{j_m}(F) \rvert
\to \lvert g_{k_n} \rvert \wedge v$ as $m \to \infty$ by continuity of
$\wedge$, so (since $\Gamma(F)$ is closed and contains all
$\lvert g_{k_n} \rvert \wedge \lvert f_{j_m}(F) \rvert$) we get
$\lvert g_{k_n} \rvert \wedge v \in \Gamma(F)$.  Letting $n \to \infty$ and
using closedness again yields $u \wedge v \in \Gamma(F)$.  Since
$u \preceq v$, we have $u \wedge v = u$, so $u \in \Gamma(F)$.

Thus $x^\pm \in \Gamma(F)$ and hence $x \in \Gamma(F)$, proving
$I(F) \subseteq \Gamma(F)$.

\medskip
\emph{Equality.}
Taking norm-closures gives
$\langle F \rangle_{\Ideal} \subseteq \Gamma(F)$.  Combined with the first
inclusion this yields $\Gamma(F) = \langle F \rangle_{\Ideal}$.

\medskip
\emph{Borel measurability.}
Fix an enumeration of all finite rational linear combinations of the vectors
$\lvert g_k \rvert \wedge \lvert f_j(F) \rvert$: for each $m \in \omega$ let
\[
    t_m(F) = \sum_{l=1}^{n(m)} q_l(m)
    \bigl(\lvert g_{k_l(m)} \rvert \wedge \lvert f_{j_l(m)}(F) \rvert\bigr),
    \qquad q_l(m) \in \bbQ.
\]
Each $t_m$ is continuous in $F \in \Subl(E)$ by
Lemma~\ref{lem:lattice-ops-lipschitz} and continuity of the selections~$f_j$.
Hence $F \mapsto (t_m(F))_{m \in \omega}$ is continuous from $\Subl(E)$ to
$E^\omega$.  By Lemma~\ref{lem:closure-map-borel} (applied to $X = E$), the
map $(x_m)_m \mapsto \overline{\{x_m \colon m \in \omega\}}$ is Borel from
$E^\omega$ to $\mathcal{F}(E)$, so
$F \mapsto \overline{\{t_m(F)\}} = \Gamma(F)$ is Borel.  Finally,
$\Ideal(E)$ is a Borel subset of $\mathcal{F}(E)$ (indeed closed
in~$\SB(E)$), so $\Gamma$ is Borel as a map into $\Ideal(E)$.
\end{proof}

\subsection{Borel generation of sublattices and ideals}\label{subsec:borel-gen}

\begin{proposition}\label{prop:borel-generated-sublattice}
Let $E$ be a separable Banach lattice.  Fix an enumeration
$(\tau_m)_{m \in \omega}$ of all \emph{lattice terms with rational coefficients}
in finitely many variables, built from variables using addition, scalar
multiplication by elements of~$\bbQ$, and the modulus operation $\lvert\cdot\rvert$.

Define $G_{\Subl} \colon E^\omega \to \Subl(E)$ by
\[
    G_{\Subl}\bigl((x_n)_{n \in \omega}\bigr)
    := \overline{\{\tau_m((x_n)_{n \in \omega}) \colon m \in \omega\}}.
\]
Then $G_{\Subl}$ is Borel measurable, and $G_{\Subl}((x_n))$ is precisely the
closed sublattice of $E$ generated by $\{x_n \colon n \in \omega\}$.

Consequently, with $\Gamma$ as in Proposition~\ref{prop:ideal-hull}, the map
\[
    G_{\Ideal} := \Gamma \circ G_{\Subl}
    \colon E^\omega \to \Ideal(E)
\]
is Borel measurable and returns the closed ideal generated by
$\{x_n \colon n \in \omega\}$.
\end{proposition}

\begin{proof}
Each term $\tau_m$ depends on finitely many coordinates and is obtained by
composing continuous operations on~$E$ (addition, rational scalar
multiplication, and modulus; see Lemma~\ref{lem:lattice-ops-lipschitz}).
Therefore the evaluation map
$E^\omega \ni (x_n) \mapsto \tau_m((x_n)) \in E$ is continuous for each~$m$,
and so
\[
    E^\omega \ni (x_n) \longmapsto
    \bigl(\tau_m((x_n))\bigr)_{m \in \omega} \in E^\omega
\]
is continuous for the product topology.

By Lemma~\ref{lem:closure-map-borel}, the closure map
$(z_m)_m \mapsto \overline{\{z_m \colon m \in \omega\}}$ is Borel from
$E^\omega$ to $\mathcal{F}(E)$, hence $G_{\Subl}$ is Borel.

Let
\[
    A((x_n)) := \{\tau_m((x_n)_{n \in \omega}) \colon m \in \omega\}.
\]
Because the family of lattice terms is closed under addition, rational scalar
multiplication, and modulus, the set $A((x_n))$ is a rational sublattice of $E$.
Since these operations are continuous on $E$
(Lemma~\ref{lem:lattice-ops-lipschitz}),
its norm-closure
\[
    \overline{A((x_n))}
    = G_{\Subl}((x_n)_{n \in \omega})
\]
is a closed sublattice of $E$.

Let $L$ be the smallest closed sublattice containing
$\{x_n \colon n \in \omega\}$.  Since each variable is a term, we have
$x_n \in \{\tau_m((x_n))\}$, so $G_{\Subl}((x_n))$ is a closed sublattice
containing $\{x_n\}$, hence $L \subseteq G_{\Subl}((x_n))$.  Conversely, any
sublattice containing $\{x_n\}$ contains the values of all lattice terms in the
$x_n$; by closedness it contains $G_{\Subl}((x_n))$.  Thus
$G_{\Subl}((x_n)) = L$.

The statement about $G_{\Ideal}$ follows immediately from
Proposition~\ref{prop:ideal-hull}.
\end{proof}

%=============================================================================
\section{Borel complexity of the Fremlin tensor product}\label{sec:tensor}
%=============================================================================

\subsection{Fremlin tensor product}

Let $E$ and $F$ be Banach lattices.  We regard the algebraic tensor product
$E \otimes F$ as canonically embedded in the Fremlin vector lattice tensor
product of $E$ and $F$, so that lattice operations such as $\lvert u \rvert$
are understood there.  The \emph{positive projective tensor norm} on
$E \otimes F$ is defined by
\[
    \lVert u \rVert_{|\pi|}
    = \inf\biggl\{\sum_{j=1}^n \lVert x_j \rVert_E \lVert y_j \rVert_F
    \colon n \in \N,\; x_j \in E^+,\; y_j \in F^+,\;
    \lvert u \rvert \preceq \sum_{j=1}^n x_j \otimes y_j\biggr\}.
\]
The completion of $(E \otimes F,\lVert\cdot\rVert_{|\pi|})$ is denoted by
$E \hat{\otimes}_{|\pi|} F$ and is called the \emph{Fremlin (projective)
tensor product} of $E$ and $F$.  It is a Banach lattice; see
\cite{Frem1974Ban} and \cite{Pug2007Fremlin}.

\subsection{Complexity of the tensor product map}

Let $\mathscr{X}$ denote the coding space $\Ideal\bigl(\FBL[\ell_1]\bigr)$,
constructed from the separable Banach lattice $E = \FBL[\ell_1]$.
Put $E_0 = E \hat{\otimes}_{|\pi|} E$, and let
$\mathscr{Y} = \Ideal(E_0)$ denote the corresponding coding space constructed
over~$E_0$.  The argument below follows the Wijsman-coordinate strategy of
Kania's tensor-kernel results for quotient-coded Banach-type structures:
fixed-coordinate distance or quotient-norm functions are first shown to be upper
semicontinuous, and Borel measurability is then read from the Wijsman subbasis;
compare \cite[Proposition~9.2 and Propositions~9.5--9.6]{Kan2026Pol}.  Here the
norm is the Fremlin positive projective tensor norm and the closedness input is
the ideal property for Banach lattices.

\begin{remark}\label{rem:ideals-not-subl}
We restrict our attention to the coding space of ideals rather than sublattices.
The Fremlin tensor product has the \emph{ideal property}: if $I \subseteq E$ and
$J \subseteq F$ are closed ideals, then $I \hat{\otimes}_{|\pi|} J$ identifies
canonically with a closed ideal of $E \hat{\otimes}_{|\pi|} F$; see, for example,
Nielsen~\cite{Nie1982Ideal} (cf.\ Fremlin~\cite{Frem1974Ban} and
Puglisi~\cite{Pug2007Fremlin}).
By contrast, the analogous statement for arbitrary closed sublattices fails in
general, so a direct tensor product map into a coding space of sublattices
of~$E_0$ is not available.
\end{remark}

\begin{lemma}\label{lem:dist-dense}
Let $(M, d)$ be a metric space and let $F \subseteq M$ be a non-empty closed set.
If $D \subseteq F$ is dense, then for every $x \in M$ one has
\[
    d(x, F) = \inf_{y \in D} d(x, y).
\]
\end{lemma}

\begin{proof}
Clearly $\inf_{y \in D} d(x, y) \ge d(x, F)$.  Conversely, given
$\varepsilon > 0$ choose $z \in F$ with $d(x, z) < d(x, F) + \varepsilon$,
and then choose $y \in D$ with $d(y, z) < \varepsilon$.  Then
$d(x, y) \le d(x, z) + \varepsilon < d(x, F) + 2\varepsilon$.
Letting $\varepsilon \downarrow 0$ yields the reverse inequality.
\end{proof}

\begin{lemma}\label{lem:dense-tensor}
Let $E$ and $F$ be Banach lattices and let $A \subseteq E$, $B \subseteq F$
be dense.  Then the linear span of $A \otimes B$ is dense in
$E \hat{\otimes}_{|\pi|} F$.
\end{lemma}

\begin{proof}
The Fremlin norm is a lattice cross norm, hence for all $x \in E$ and $y \in F$,
\begin{equation}\label{eq:crossnorm-est}
    \lVert x \otimes y \rVert_{|\pi|}
    \le \lVert x \rVert_E \, \lVert y \rVert_F;
\end{equation}
see, e.g., Fremlin~\cite{Frem1974Ban} (or Puglisi~\cite{Pug2007Fremlin}).
Consequently, the canonical bilinear map $(x, y) \mapsto x \otimes y$ is jointly
continuous.

Let $u = \sum_{l=1}^k x_l \otimes y_l \in E \otimes F$ and fix
$\varepsilon > 0$.  Choose $a_l \in A$ and $b_l \in B$ such that
\[
    \lVert x_l - a_l \rVert_E
    < \min\Bigl\{1,\;\frac{\varepsilon}{4k(\lVert y_l \rVert_F + 1)}\Bigr\},
    \qquad
    \lVert y_l - b_l \rVert_F
    < \frac{\varepsilon}{4k(\lVert x_l \rVert_E + 1)}.
\]
Then $\lVert a_l \rVert_E \le \lVert x_l \rVert_E + 1$.  Using
$x_l \otimes y_l - a_l \otimes b_l
= (x_l - a_l) \otimes y_l + a_l \otimes (y_l - b_l)$
and~\eqref{eq:crossnorm-est}, we obtain
\[
\begin{aligned}
    \lVert x_l \otimes y_l - a_l \otimes b_l \rVert_{|\pi|}
    &\le \lVert x_l - a_l \rVert_E \, \lVert y_l \rVert_F
       + \lVert a_l \rVert_E \, \lVert y_l - b_l \rVert_F \\
    &< \frac{\varepsilon}{4k} + \frac{\varepsilon}{4k}
    = \frac{\varepsilon}{2k}.
\end{aligned}
\]
Summing over $l = 1, \ldots, k$ gives
\[
    \biggl\lVert u - \sum_{l=1}^k a_l \otimes b_l \biggr\rVert_{|\pi|}
    < \varepsilon.
\]
Since finite sums are dense in the completion $E \hat{\otimes}_{|\pi|} F$, the
linear span of $A \otimes B$ is dense in $E \hat{\otimes}_{|\pi|} F$.
\end{proof}

\begin{theorem}\label{thm:tensor}
The map
\[
    \Xi \colon \mathscr{X} \times \mathscr{X} \ni (J_1, J_2)
    \longmapsto J_1 \hat{\otimes}_{|\pi|} J_2 \in \mathscr{Y}
\]
is $\boldsymbol{\Sigma}^0_2$-measurable when $\mathscr{X}$ and $\mathscr{Y}$
carry the (subspace) Wijsman topologies.  In particular, $\Xi$ is Borel
measurable.  Moreover, $\Graph(\Xi)$ is a $G_\delta$ subset of
$\mathscr{X} \times \mathscr{X} \times \mathscr{Y}$.
\end{theorem}

\begin{proof}
Fix a compatible metric $d$ on $E_0$ and a dense sequence
$(\alpha_n)_{n \in \omega}$ in~$E_0$.  Recall that the Wijsman topology on
$\mathcal{F}(E_0)$ is the initial topology for the distance maps
$F \mapsto d(\alpha_n, F)$, $n \in \omega$.

\medskip
\emph{A continuous dense sequence of simple tensors.}
Let $f_j \colon \mathscr{X} \to E$, $j \in \omega$, be the continuous selections
from Theorem~\ref{thm:selection}.  For
$(J_1, J_2) \in \mathscr{X} \times \mathscr{X}$ and each index $m \in \omega$
coding a tuple $(k, q_1, \ldots, q_k, (i_1, j_1), \ldots, (i_k, j_k))$ with
$q_l \in \bbQ$, set
\[
    s_m(J_1, J_2)
    = \sum_{l=1}^k q_l \bigl(f_{i_l}(J_1) \otimes f_{j_l}(J_2)\bigr) \in E_0.
\]
Since the bilinear map $(x, y) \mapsto x \otimes y \colon E \times E \to E_0$ is
continuous (as a bounded bilinear map between Banach spaces) and addition and
scalar multiplication are continuous, each $s_m$ is continuous.  The set
$\{s_m(J_1, J_2) \colon m \in \omega\}$ is dense in
$J_1 \hat{\otimes}_{|\pi|} J_2$.  Indeed,
$\{f_i(J_1) \colon i \in \omega\}$ is dense in $J_1$ and
$\{f_j(J_2) \colon j \in \omega\}$ is dense in $J_2$ by
Theorem~\ref{thm:selection}, so Lemma~\ref{lem:dense-tensor} yields density of
the span of $\{f_i(J_1) \otimes f_j(J_2)\}$, and passing to rational
coefficients does not change the closure.

\medskip
\emph{Wijsman coordinates are upper semicontinuous.}
For each $n \in \omega$ consider
\[
    \varphi_n(J_1, J_2)
    := d\bigl(\alpha_n,\, J_1 \hat{\otimes}_{|\pi|} J_2\bigr)
    = \inf_{m \in \omega} d\bigl(\alpha_n,\, s_m(J_1, J_2)\bigr),
\]
by Lemma~\ref{lem:dist-dense}.  For each fixed $m$, the map
$(J_1, J_2) \mapsto d(\alpha_n, s_m(J_1, J_2))$ is continuous.  Hence $\varphi_n$
is the pointwise infimum of countably many continuous functions, and therefore
$\varphi_n$ is upper semicontinuous.

\medskip
\emph{$\boldsymbol{\Sigma}^0_2$-measurability of $\Xi$.}
A subbasic open set in the Wijsman topology on $\mathscr{Y}$ is of the form
\[
    U(n,r)^{<} = \{K \in \mathscr{Y} \colon d(\alpha_n, K) < r\}
    \quad\text{or}\quad
    U(n,r)^{>} = \{K \in \mathscr{Y} \colon d(\alpha_n, K) > r\},
\]
where $n \in \omega$ and $r \in \R$.  We have
\[
    \Xi^{-1}\bigl(U(n,r)^{<}\bigr)
    = \{(J_1, J_2) \colon \varphi_n(J_1, J_2) < r\},
\]
which is open because $\varphi_n$ is upper semicontinuous.  Moreover,
\[
    \Xi^{-1}\bigl(U(n,r)^{>}\bigr)
    = \{(J_1, J_2) \colon \varphi_n(J_1, J_2) > r\}
    = \bigcup_{\substack{q \in \bbQ \\ q > r}}
      \{(J_1, J_2) \colon \varphi_n(J_1, J_2) \ge q\},
\]
and each set $\{\varphi_n \ge q\}$ is closed since $\varphi_n$ is upper
semicontinuous.  Hence $\Xi^{-1}(U(n,r)^{>})$ is $F_\sigma$.
Since $\mathscr{X} \times \mathscr{X}$ is metrisable, every open subset of it is
also $F_\sigma$. Therefore the preimage of each subbasic Wijsman-open set is
$F_\sigma$, and so the preimage of every open subset of $\mathscr{Y}$ (being a
union of finite intersections of subbasic opens) is $F_\sigma$.
Thus $\Xi$ is $\boldsymbol{\Sigma}^0_2$-measurable.

\medskip
\emph{The graph is $G_\delta$.}
Fix $n \in \omega$ and define on
$\mathscr{X} \times \mathscr{X} \times \mathscr{Y}$ the functions
\[
    a_n(J_1, J_2, K) = d(\alpha_n, K), \qquad
    b_n(J_1, J_2, K) = \varphi_n(J_1, J_2).
\]
Then $a_n$ is continuous and $b_n$ is upper semicontinuous.  Let
\[
    A_n = \{(J_1, J_2, K) \colon a_n(J_1, J_2, K) = b_n(J_1, J_2, K)\}.
\]
Write $A_n = \{a_n \le b_n\} \cap \{b_n \le a_n\}$.
The set $\{a_n \le b_n\}$ is closed because its complement
$\{a_n > b_n\} = \bigcup_{q \in \bbQ} \{a_n > q\} \cap \{b_n < q\}$ is open.
Also, $\{b_n \le a_n\}$ is $G_\delta$ because its complement
\[
    \{b_n > a_n\}
    = \bigcup_{q \in \bbQ} \{a_n < q\} \cap \{b_n > q\}
\]
is $F_\sigma$: the sets $\{a_n < q\}$ are open, hence $F_\sigma$ in the
metrisable space $\mathscr{X} \times \mathscr{X} \times \mathscr{Y}$, while the
sets $\{b_n > q\}$ are $F_\sigma$ because $b_n$ is upper semicontinuous.
Therefore $A_n$ is $G_\delta$.  Finally,
\[
    \Graph(\Xi) = \bigcap_{n \in \omega} A_n,
\]
so $\Graph(\Xi)$ is $G_\delta$.
\end{proof}

\begin{corollary}\label{cor:tensor-sections}
Fix $J_1 \in \mathscr{X}$.  The section
\[
    \Xi_{J_1} \colon \mathscr{X} \to \mathscr{Y}, \qquad
    \Xi_{J_1}(J_2) = \Xi(J_1, J_2),
\]
is $\boldsymbol{\Sigma}^0_2$-measurable.  In particular, it is Borel measurable,
and it is continuous on a dense $G_\delta$ subset of $\mathscr{X}$.
The analogous statement holds for the sections $J_1 \mapsto \Xi(J_1, J_2)$
with $J_2$ fixed.
\end{corollary}

\begin{proof}
The $\boldsymbol{\Sigma}^0_2$-measurability follows by restricting
Theorem~\ref{thm:tensor} to $\{J_1\} \times \mathscr{X}$.
Since both $\mathscr{X}$ and $\mathscr{Y}$ are Polish, every
$\boldsymbol{\Sigma}^0_2$-measurable map from $\mathscr{X}$ to
$\mathscr{Y}$ is of Baire class~$1$.
The statement about continuity on a dense $G_\delta$ set is then the
classical theorem that every Baire class~$1$ map from a Polish space
into a metric space has a comeagre set of points of continuity.
The analogous statement for the other sections is identical.
\end{proof}

\begin{remark}\label{rem:closure-not-cont}
It is tempting to factor $\Xi$ as the composition of the continuous map
$S \colon \mathscr{X} \times \mathscr{X} \to E_0^\omega$,
$S(J_1, J_2) = \bigl(s_m(J_1, J_2)\bigr)_{m \in \omega}$, with the closure map
\[
    E_0^\omega \ni (z_j)_{j \in \omega}
    \longmapsto \overline{\{z_j \colon j \in \omega\}} \in \mathcal{F}(E_0),
\]
and to deduce continuity of $\Xi$.  However, this closure map is
\emph{not} continuous for the product topology on $E_0^\omega$ and the Wijsman
topology on $\mathcal{F}(E_0)$.  Indeed, consider $E_0 = \R$ and define
$z^{(n)} \in \R^\omega$ by $z^{(n)}_n = -1$ and $z^{(n)}_j = 0$ for $j \ne n$.
Then $z^{(n)} \to 0^\omega$ in the product topology, but
$d\bigl(-1, \overline{\{z^{(n)}_j\}}\bigr) = 0$ for all~$n$, whereas
$d\bigl(-1, \overline{\{0, 0, \ldots\}}\bigr) = 1$.  Thus continuity of~$\Xi$
cannot be established by this factorisation.
\end{remark}

%=============================================================================
%\section{Towards equivalence of coding spaces}\label{sec:equiv}
%=============================================================================

%We now discuss the relationship between the two coding spaces,
% $\Subl(\mathcal{C})$ and $\Ideal(\FBL[\ell_1])$, following Bossard's
% rigidification philosophy~\cite{Boss2002Cod} and Kania's quotient-generator
% viewpoint~\cite[Main Theorem~3.13]{Kan2026Pol}.

% \subsection{A Borel normalisation lemma}

% The first ingredient is a way to extract, from the continuous selections of
% Theorem~\ref{thm:selection}, a Borel-measurable dense sequence in the unit sphere
% of each sublattice.

% \begin{lemma}\label{lem:borel-unit-sphere-seq}
% Let $X$ be a separable Banach space and let $f_j \colon \SB(X) \to X$,
% $j \in \omega$, be continuous selections as in Theorem~\ref{thm:selection}.
% Define $\nu \colon X \to X$ by
% \[
%     \nu(x) =
%     \begin{cases}
%     x / \lVert x \rVert_X, & x \ne 0, \\
%     0, & x = 0.
%     \end{cases}
% \]
% Then the map $\iota \colon \SB(X) \to X^\omega$ given by
% $\iota(F) = \bigl(\nu(f_j(F))\bigr)_{j \in \omega}$ is Borel measurable.
% Moreover, for every $F \in \SB(X)$ with $F \ne \{0\}$, the set
% $\{\nu(f_j(F)) \colon f_j(F) \ne 0\}$ is dense in the unit sphere
% $\mathscr{S}_F$.
% \end{lemma}

% \begin{proof}
% For each fixed $j$, the map $F \mapsto f_j(F)$ is continuous.
% It therefore suffices to note that $\nu$ is Borel measurable on~$X$.
% Indeed, $\nu$ is continuous on $X \setminus \{0\}$ and
% $\nu^{-1}(\{0\}) = \{0\}$.  Hence for every open $U \subseteq X$,
% \[
%     \nu^{-1}(U)
%     = \Bigl[(X \setminus \{0\}) \cap (\nu|_{X \setminus \{0\}})^{-1}(U)\Bigr]
%     \;\cup\;
%     \begin{cases}
%     \{0\}, & 0 \in U, \\
%     \emptyset, & 0 \notin U,
%     \end{cases}
% \]
% which is Borel in~$X$.  Thus $\nu$ is Borel, and so $F \mapsto \nu(f_j(F))$ is
% Borel measurable as a composition of a continuous map with a Borel map.
% Therefore $\iota$ is Borel as a map into the product space $X^\omega$.

% If $F \ne \{0\}$ and $x \in \mathscr{S}_F$, pick indices $(j_n)_{n \in \omega}$
% such that $f_{j_n}(F) \to x$.  Then
% $\lVert f_{j_n}(F) \rVert \to \lVert x \rVert = 1$, so $f_{j_n}(F) \ne 0$
% eventually, and $\nu(f_{j_n}(F)) \to \nu(x) = x$ by continuity of $\nu$ at
% $x \ne 0$.
% \end{proof}

% \subsection{The quotient map construction}

% Recall that $\mathcal{C} = C(\Delta; L_1)$ and write
% $\overline B_{\mathcal C}$ for its closed unit ball.  Given a sequence
% $\tilde{f} = (f_k)_{k \in \omega} \in \overline B_{\mathcal C}^{\omega}$,
% define first the contractive linear map
% $T_{\tilde f}\colon \ell_1\to\mathcal C$ by
% $T_{\tilde f}(e_k)=f_k$.  The universal property of $\FBL[\ell_1]$ then gives
% a unique contractive lattice homomorphism
% $\mathfrak{Q}_{\tilde{f}} \colon \FBL[\ell_1] \to \mathcal{C}$ satisfying
% $\mathfrak{Q}_{\tilde{f}}(\delta_{e_k}) = f_k$ on the generators.

% When $\tilde{f} = \iota(Y)$ for some $Y \in \Subl(\mathcal{C})$ with
% $Y \ne \{0\}$, the image of $\mathfrak{Q}_{\tilde{f}}$ contains the dense
% subset $\{\nu(f_k(Y)) \colon f_k(Y) \ne 0\}$ of~$\mathscr{S}_Y$.

% \begin{remark}\label{rem:range-not-closed}
% In general, a contractive lattice homomorphism between Banach lattices need not
% have closed range.  For instance, the natural inclusion
% $\ell_1 \hookrightarrow \ell_2$ is a contractive lattice homomorphism
% (with the coordinatewise order) whose range is not closed in~$\ell_2$.
% Consequently, the map $\mathfrak{Q}_{\tilde{f}}$ constructed above may fail to
% be surjective onto~$Y$, and the assertion that its image equals~$Y$ requires
% additional justification.
% \end{remark}
% 
\subsection{Open problems}\label{subsec:open}

% In view of Remark~\ref{rem:range-not-closed}, we formulate the desired
% equivalence between coding spaces as an open problem.

% \begin{problem}\label{prob:coding-equiv}
% Does there exist an injective Borel-measurable map
% \[
%     \Psi \colon \Subl(\mathcal{C})
%     \longrightarrow \Ideal\bigl(\FBL[\ell_1]\bigr)
% \]
% such that, for each $Y \in \Subl(\mathcal{C})$, the
% quotient $\FBL[\ell_1] / \Psi(Y)$ is isometrically lattice-isomorphic to~$Y$?
% \end{problem}

% A natural candidate is $\Psi(Y) = \ker \mathfrak{Q}_{\iota(Y)}$, where $\iota$
% is the Borel normalisation from Lemma~\ref{lem:borel-unit-sphere-seq}.
% Establishing this requires resolving two difficulties: ensuring that the
% quotient map $\mathfrak{Q}_{\iota(Y)}$ is surjective onto~$Y$ (or, equivalently,
% has closed range equal to~$Y$), and proving that $\Psi$ is injective as a map
% between coding spaces (not merely up to isometric lattice isomorphism).

% A positive resolution would likely require a Bossard-style rigidification
% argument~\cite{Boss2002Cod} that pins down a canonical copy of the quotient
% inside~$\mathcal{C}$.  We also note that the Borel measurability of the kernel
% map
% $\tilde{f} \mapsto \ker \mathfrak{Q}_{\tilde{f}}$ is a non-trivial question that
% would benefit from a systematic study of the space of contractive lattice
% homomorphisms $\FBL[\ell_1] \to \mathcal{C}$ as a Polish space.

% Kania's framework provides faithful quotient-presentation codes and a general
% definability scheme for such quotient spaces
% \cite[Main Theorem~3.13 and Theorem~3.16]{Kan2026Pol}.  Problem~\ref{prob:coding-equiv}
% asks for a stronger bridge: a Borel way of passing from the embedding code
% \(Y\subseteq\mathcal C\) to a quotient kernel while retaining the chosen copy of
% \(Y\), not only its abstract isometric lattice-isomorphism class.

%=============================================================================
\section{Further directions}\label{sec:further}
%=============================================================================

We briefly mention several classes of Banach lattices whose descriptive
complexity within the coding spaces introduced above would be natural to study.
For properties expressible by quotient-norm formulae over a fixed dense set,
Kania's definability theorem gives a general source of Borel upper bounds
\cite[Theorem~3.16]{Kan2026Pol}; the questions below indicate where additional
Banach-lattice structure or the sublattice coding creates extra work.

\subsection{Lattice properties of interest}

\begin{definition}
Let $E$ be a Banach lattice.
\begin{enumerate}[label=(\roman*)]
    \item $E$ has the \emph{Fatou property} if, for every upward-directed set
      $A \subseteq E^+$ with supremum~$b$, one has
      $\lVert b \rVert_E = \sup\{\lVert a \rVert_E \colon a \in A\}$.
    \item $E$ has the \emph{Nakano property} if, for every upward-directed and
      order-bounded set $A \subseteq E^+$,
      $\inf\{\lVert b \rVert_E \colon b \text{ is an upper bound of } A\}
      = \sup\{\lVert a \rVert_E \colon a \in A\}$.
    \item $E$ has the \emph{strong Nakano property} if, for every
      upward-directed and norm-bounded set $A \subseteq E^+$, there exists an
      upper bound $b \in E^+$ with
      $\lVert b \rVert_E = \sup\{\lVert a \rVert_E \colon a \in A\}$.
\end{enumerate}
\end{definition}

By Wickstead~\cite{Wick2007nak}, a Banach lattice $E$ has the Nakano property if
and only if its Dedekind completion $\hat{E}$ has the Fatou property.  By
Avil\'es et al.~\cite{aviles2018free}, the free Banach lattice generated by any
set has the strong Nakano property.

\begin{definition}\label{def:semiprojective}
A separable Banach lattice $E$ is \emph{semiprojective} if, for every Banach lattice $F$,
every increasing sequence
$(J_n)_{n=1}^\infty$ of closed ideals in~$F$ with
$J = \overline{\bigcup_{n=1}^\infty J_n}$, every lattice homomorphism
$\phi \colon E \to F/J$, and every $\varepsilon > 0$, there exist an integer
$n \ge 1$ and a lattice homomorphism
$\psi \colon E \to F/J_n$ such that
$\pi_n \circ \psi = \phi$ and
$\lVert \psi \rVert \le (1 + \varepsilon) \lVert \phi \rVert$,
where $\pi_n \colon F/J_n \to F/J$ denotes the canonical quotient map.
A lattice homomorphism $\phi$ for which such a $\psi$ exists is said to be
\emph{partially liftable}.
\end{definition}

\begin{definition}\label{def:projective}
Let $F$ be a Banach lattice, $J \subseteq F$ a closed ideal, and
$Q \colon F \to F/J$ the canonical quotient map.  A lattice homomorphism
$\phi \colon E \to F/J$ is \emph{liftable} if there exists a lattice
homomorphism $\psi \colon E \to F$ such that $Q \circ \psi = \phi$ and
$\lVert \psi \rVert \le (1 + \varepsilon) \lVert \phi \rVert$ for a given
$\varepsilon > 0$:
\[
    \begin{tikzcd}
    E \arrow[r, "\psi"] \arrow[rd, "\phi"'] & F \arrow[d, "Q"] \\
                                              & F/J
    \end{tikzcd}
\]
The Banach lattice $E$ is \emph{projective} if, for
every Banach lattice~$F$, every closed ideal $J \subseteq F$, and every
$\varepsilon > 0$, every lattice homomorphism $\phi \colon E \to F/J$ is
liftable (with the given~$\varepsilon$).
\end{definition}

As noted in~\cite{de2015free}, one cannot in general take $\varepsilon = 0$, even
when $E \cong \R$; the approximate lifting condition
$\lVert \psi \rVert \le (1 + \varepsilon) \lVert \phi \rVert$ is therefore
essential.

\begin{remark}\label{rem:proj-examples}
De~Pagter and Wickstead prove that free Banach lattices generated by sets
are projective.  They also prove separately that every finite-dimensional
Banach lattice is projective.  In particular, $\R$ is projective.

This gives a basic example of a projective Banach lattice that is not free:
the free Banach lattice generated by a singleton is lattice-isometric to
$\ell_\infty^2$, not to~$\R$.
\end{remark}

\begin{remark}\label{rem:proj-structure}
A permanence property that is available in the literature is the following:
if $(P_n)_{n=1}^\infty$ are projective Banach lattices with topological order
units (i.e.\ quasi-interior points: positive elements whose principal ideal is
norm-dense in the whole space), then their $\ell_1$-sum is projective
\cite[Theorem~11.8]{de2015free}.

By contrast, the corresponding finite $\ell_\infty$-sum question is explicitly
left open in \cite[Question~12.12]{de2015free}.  Thus we do not use any
closure property for finite $\ell_\infty$-direct sums here.
\end{remark}

\subsection{Principal ideals and metrisable SICK compacta}\label{subsec:principal-ideals}

Let $E$ be a Banach lattice and $u \in E$ a non-zero positive element.  The
principal ideal $J_u$ need not be closed in the ambient norm of~$E$ (for
instance, taking $E = L_1[0,1]$ and $u = \mathbf{1}$ gives $J_u = L_\infty$,
which is not norm-closed in~$L_1$).  However, equipped with the \emph{gauge norm}
\[
    \lVert x \rVert_u := \inf\{r > 0 \colon \lvert x \rvert \preceq r u\},
\]
the ideal $J_u$ becomes an AM-space with unit~$u$, hence by the
Krein--Kakutani representation theorem is lattice-isomorphic to $C(K; \R)$ for
some compact Hausdorff space~$K$.  When the ambient norm-closure is needed we
write $\overline{J_u}^{\lVert\cdot\rVert_E}$.

By Schaefer~\cite{Scha1974Ban}, $K$ can be taken to be the space $K_u(E)$ of
all valuations of $E$ preserving~$u$, with the topology of pointwise convergence.
Here, a valuation preserving~$u$ is a map $\varphi \colon E \to [0, \infty]$ such
that $\varphi(u) = 1$, $\varphi(x + y) = \varphi(x) + \varphi(y)$ and
$\varphi(x \wedge y) = \varphi(x) \wedge \varphi(y)$ for all $x, y \in E^+$,
and $\varphi(r z) = \lvert r \rvert \varphi(\lvert z \rvert)$ for all
$r \in \R$ and $z \in E$ (with the convention $0 \cdot \infty = 0$).

\begin{definition}[Metrisable SICK compacta; cf.~{\cite{AMRT2022Sick}}]
A compact metrisable space $K$ is a \emph{metrisable SICK compactum} (Separable
Ideal $C(K)$) if $K$ is homeomorphic to $K_u(E)$ for some separable Banach
lattice $E$ and some non-zero $u \in E^+$.
\end{definition}

Avil\'es, Mart\'\i nez-Cervantes, Rueda Zoca, and Tradacete define SICK compacta
in the class of all compact Hausdorff spaces.  In this paper we restrict
explicitly to the metrisable part of that class.  This restriction is important
for descriptive set theory: compact metrisable spaces have a standard
hyperspace coding, for instance as closed subsets of the Hilbert cube
$[0,1]^\omega$, whereas non-metrisable compact Hausdorff spaces require a
different coding.  We do not use the non-metrisable part of the SICK theory
below.

\subsection{Monotonicity and order smoothness in the sense of Kurc}\label{subsec:kurc}

We now turn to several lattice-geometric properties introduced by
Kurc~\cite{Kurc1993Dual} and determine their descriptive complexity within
$\Subl(E)$.

\begin{definition}[Kurc~\cite{Kurc1993Dual}]\label{def:kurc}
Let $X$ be a Banach lattice.
\begin{enumerate}[label=(\roman*)]
    \item $X$ is \emph{strictly monotone} (STM) if $\lVert x - y \rVert < \lVert x \rVert$
      whenever $\mathbf{0}_X \prec y \preceq x$.
    \item For $\varepsilon \in (0, 1]$ define
      \[
          \delta_X(\varepsilon)
          = \inf\bigl\{1 - \lVert x - y \rVert \colon
            \mathbf{0}_X \preceq y \preceq x,\;
            \lVert x \rVert = 1,\;
            \lVert y \rVert \ge \varepsilon\bigr\}.
      \]
      Then $X$ is \emph{uniformly monotone} (UM) if $\delta_X(\varepsilon) > 0$
      for all $\varepsilon \in (0, 1]$.
    \item $X$ is \emph{order smooth}, in abbreviation \emph{(o)-Sm}, if for
      every $x$ in the positive part of the unit sphere $S(X_+)$, the set
      $\partial_+\lVert x \rVert
      = \{x^* \in S(X^*_+) \colon \langle x, x^* \rangle = \lVert x \rVert\}$
      contains no proper order interval.
    \item For $\tau \in [0, 1]$ define the \emph{order modulus of smoothness}
      \[
          \rho_X(\tau)
          = \sup\bigl\{\lVert x \vee \tau y \rVert - 1 \colon
            \mathbf{0}_X \preceq x, y,\;
            \lVert x \rVert = \lVert y \rVert = 1\bigr\}.
      \]
      Then $X$ is \emph{order uniformly smooth}, in abbreviation
      \emph{(o)-USm}, if $\rho_X(\tau) / \tau \to 0$ as $\tau \downarrow 0$.
\end{enumerate}
\end{definition}

Kurc proves the identity
$\delta_X(\varepsilon) = \sigma_X(\varepsilon) / (1 + \sigma_X(\varepsilon))$,
where
$\sigma_X(\varepsilon)
= \inf\{\lVert x + y \rVert - 1 \colon
  x, y \succeq \mathbf{0}_X,\;
  \lVert x \rVert = 1,\;
  \lVert y \rVert \ge \varepsilon\}$,
so that $\delta_X(\varepsilon) > 0$ if and only if
$\sigma_X(\varepsilon) > 0$; see~\cite{Kurc1993Dual}.  He also establishes a
duality between UM and (o)-USm: if $X$ is UM then $X^*$ is (o)-USm, and
conversely.

For later use, when $F = \{0\}$ we adopt the conventions
\[
    \sigma_{\{0\}}(\varepsilon) = 0,
    \qquad
    \rho_{\{0\}}(\tau) = 0
    \qquad (\varepsilon,\tau \in (0,1]).
\]
This convention treats the zero lattice uniformly in the formulae below.  If one
prefers to regard the zero lattice as uniformly monotone by vacuity, the
corresponding class differs from ours by the singleton $\{\{0\}\}$, which does
not affect any of the Borel upper bounds.

\begin{theorem}\label{thm:kurc-borel}
Let $E$ be a separable Banach lattice and endow $\Subl(E)$ with the Wijsman
topology.
\begin{enumerate}[label=\textup{(\alph*)}]
    \item\label{kurc:baire1}
      For each fixed rational $\varepsilon, \tau \in (0, 1]$, the maps
      $F \mapsto \sigma_F(\varepsilon)$ and $F \mapsto \rho_F(\tau)$ are of
      Baire class~$1$ on $\Subl(E)$.  Consequently, the sets
      $\{F \colon \sigma_F(\varepsilon) > 0\}$ and
      $\{F \colon \rho_F(\tau) < r\}$ (for rational~$r$) are $F_\sigma$ subsets
      of $\Subl(E)$.
    \item\label{kurc:UM}
      The class $\mathrm{UM}(E)$ of uniformly monotone lattices in $\Subl(E)$
      is a $\boldsymbol{\Pi}^0_3$ subset of $\Subl(E)$.
    \item\label{kurc:oUSm}
      The class $\mathrm{(o)\text{-}USm}(E)$ of order uniformly smooth lattices
      in $\Subl(E)$ is a $\boldsymbol{\Pi}^0_3$ subset of $\Subl(E)$.
    \item\label{kurc:STM}
      The class $\mathrm{STM}(E)$ of strictly monotone lattices in $\Subl(E)$
      is coanalytic ($\boldsymbol{\Pi}^1_1$).
\end{enumerate}
\end{theorem}

\begin{proof}
Let $f_j \colon \SB(E) \to E$, $j \in \omega$, be continuous selections as in
Theorem~\ref{thm:selection}.

\medskip
\emph{Proof of~\ref{kurc:baire1}.}
Fix a rational $\varepsilon \in (0, 1]$.

\smallskip
\emph{Reduction of $\sigma_F(\varepsilon)$ to unit vectors.}
Since the norm is monotone on $F_+$, for $x \in S(F_+)$ and $y \in F_+$ with
$\lVert y \rVert \ge \varepsilon$ we have
$0 \preceq \varepsilon\, y / \lVert y \rVert \preceq y$, hence
\[
    \lVert x + \varepsilon\, y / \lVert y \rVert \rVert
    \le \lVert x + y \rVert.
\]
Therefore, for every $F \ne \{0\}$, the infimum in the definition of
$\sigma_F(\varepsilon)$ may be
computed already on vectors with $\lVert y \rVert = \varepsilon$, and
\begin{equation}\label{eq:sigma-unit-reduction}
    \sigma_F(\varepsilon)
    = \inf\bigl\{\lVert x + \varepsilon\, y \rVert - 1
      \colon x, y \in S(F_+)\bigr\}.
\end{equation}

\smallskip
\emph{Continuous truncated normalisations.}
For each $n \ge 1$ define
\[
    \nu_n(x) = \frac{x}{\max\{\lVert x \rVert,\, 1/n\}}, \qquad x \in E.
\]
Then $\nu_n$ is continuous on $E$: writing
$\nu_n(x) = x / h_n(\lVert x \rVert)$ where $h_n(t) = \max\{t, 1/n\}$, the
denominator is a continuous function bounded below by $1/n > 0$.  Moreover,
$\lVert \nu_n(x) \rVert \le 1$ for all $x$, and
$\nu_n(x) = x / \lVert x \rVert$ whenever $\lVert x \rVert \ge 1/n$.

Set
\[
    u_{j,n}(F) := \nu_n\bigl(\lvert f_j(F) \rvert\bigr) \in E, \qquad
    F \in \Subl(E).
\]
Each map $F \mapsto u_{j,n}(F)$ is continuous on $\Subl(E)$: the selections
$f_j$ are continuous, the modulus
$x \mapsto \lvert x \rvert$ is continuous
(Lemma~\ref{lem:lattice-ops-lipschitz}\,(i)), and $\nu_n$ is continuous.

We record three properties of the truncated selectors.
Write
\[
    u_{j,n}(F)=c_{j,n}(F)\,\lvert f_j(F)\rvert,
    \qquad
    c_{j,n}(F):=\frac{1}{\max\{\lVert f_j(F)\rVert,\,1/n\}}.
\]
Then $0 \preceq u_{j,n}(F)\in F$, $\lVert u_{j,n}(F)\rVert\le 1$, and for each
fixed $j$ and $F$ the scalars $c_{j,n}(F)$ are non-decreasing in $n$.

Moreover, if $F\ne\{0\}$ and $x\in S(F_+)$, choose indices $(j_m)_{m\in\omega}$
such that $f_{j_m}(F)\to x$. Since $\lVert f_{j_m}(F)\rVert\to 1$, for every fixed
$n$ we eventually have $\lVert f_{j_m}(F)\rVert\ge 1/n$, and hence
\[
    u_{j_m,n}(F)
    = \frac{\lvert f_{j_m}(F)\rvert}{\lVert f_{j_m}(F)\rVert}
    \longrightarrow x.
\]
Thus, for each fixed $n$, the set
\[
    \{u_{j,n}(F)\colon j\in\omega,\ \lVert f_j(F)\rVert\ge 1/n\}
\]
is dense in $S(F_+)$.

\smallskip
\emph{Upper semicontinuous approximants for $\sigma_F(\varepsilon)$.}
For $i, j \in \omega$ and $n \ge 1$ define
\[
    \psi_{i,j,n}(F)
    := \lVert u_{i,n}(F) + \varepsilon\, u_{j,n}(F) \rVert - 1
    + 2\bigl(1 - \lVert u_{i,n}(F) \rVert\bigr)
    + 2\bigl(1 - \lVert u_{j,n}(F) \rVert\bigr).
\]
Each $\psi_{i,j,n}$ is a \emph{continuous} function on $\Subl(E)$ (a
composition and sum of continuous maps).  The penalty terms
$2(1 - \lVert u_{i,n}(F) \rVert)$ are non-negative and vanish precisely when
$\lVert u_{i,n}(F) \rVert = 1$; when $u_{i,n}(F) = 0$ (i.e.\ $f_i(F) = 0$),
the penalty contributes~$2$, ensuring such indices cannot artificially decrease
the infimum.

Define
\[
    \sigma^{(n)}_F(\varepsilon)
    := \inf_{i,j \in \omega} \psi_{i,j,n}(F), \qquad F \in \Subl(E).
\]
Since $\sigma^{(n)}_\bullet(\varepsilon)$ is the pointwise infimum of countably
many \emph{continuous} functions, it is upper semicontinuous, hence Baire
class~$1$.

We claim that $\sigma^{(n)}_F(\varepsilon) \downarrow \sigma_F(\varepsilon)$
for every $F \ne \{0\}$.

Fix $F \ne \{0\}$.  First note that for every fixed $i,j$ the sequence
$\psi_{i,j,n}(F)$ is non-increasing in $n$.
Indeed, write
\[
    u_{i,n}(F) = c_{i,n}\,\lvert f_i(F) \rvert,
    \qquad
    u_{j,n}(F) = c_{j,n}\,\lvert f_j(F) \rvert,
\]
where $0 \le c_{i,n} \le c_{i,n+1}$ and
$0 \le c_{j,n} \le c_{j,n+1}$.
If only the first coefficient is increased from $c$ to $c' \ge c$, while
the second vector is kept fixed, then by the triangle inequality the main norm
term can increase by at most
$(c'-c)\,\lVert f_i(F) \rVert$, whereas the corresponding penalty term
decreases by exactly
$2(c'-c)\,\lVert f_i(F) \rVert$.
Hence the total expression does not increase.
The same argument applies to the second coefficient, since
$\varepsilon \le 1 < 2$.
Therefore $\psi_{i,j,n+1}(F) \le \psi_{i,j,n}(F)$, and taking infima gives that
$\sigma^{(n)}_F(\varepsilon)$ is non-increasing.

Next we prove that
$\sigma_F(\varepsilon) \le \sigma^{(n)}_F(\varepsilon)$ for every $n$.
Fix $i,j,n$ and put
\[
    a := u_{i,n}(F), \qquad b := u_{j,n}(F), \qquad
    \alpha := \lVert a \rVert, \qquad \beta := \lVert b \rVert.
\]
If $\alpha,\beta > 0$, then
$x := a/\alpha$ and $y := b/\beta$ belong to $S(F_+)$, and
\[
    x + \varepsilon y
    = (a + \varepsilon b)
      + (1-\alpha)x
      + \varepsilon(1-\beta)y.
\]
Hence
\[
    \lVert x + \varepsilon y \rVert - 1
    \le \lVert a + \varepsilon b \rVert - 1
      + (1-\alpha) + \varepsilon(1-\beta)
    \le \psi_{i,j,n}(F),
\]
because $\varepsilon \le 1$.
Using~\eqref{eq:sigma-unit-reduction}, we obtain
$\sigma_F(\varepsilon) \le \psi_{i,j,n}(F)$.
If $\alpha=0$ or $\beta=0$, then
$\psi_{i,j,n}(F) \ge \varepsilon \ge \sigma_F(\varepsilon)$, since
$\sigma_F(\varepsilon) \le \varepsilon$ by the triangle inequality.
Taking the infimum over $i,j$ yields
\[
    \sigma_F(\varepsilon) \le \sigma^{(n)}_F(\varepsilon).
\]

Conversely, let $\eta > 0$.
By~\eqref{eq:sigma-unit-reduction}, choose $x,y \in S(F_+)$ such that
\[
    \lVert x + \varepsilon y \rVert - 1
    < \sigma_F(\varepsilon) + \eta.
\]
Since $\{f_j(F)\}_{j \in \omega}$ is dense in $F$ and the map
$z \mapsto \lvert z \rvert/\lVert z \rVert$ is continuous at every non-zero
positive vector, we may choose indices $i,j$ such that
\[
    \left\lVert
    \frac{\lvert f_i(F) \rvert}{\lVert f_i(F) \rVert} - x
    \right\rVert < \eta,
    \qquad
    \left\lVert
    \frac{\lvert f_j(F) \rvert}{\lVert f_j(F) \rVert} - y
    \right\rVert < \eta.
\]
For all sufficiently large $n$, we have
$\lVert f_i(F) \rVert,\lVert f_j(F) \rVert \ge 1/n$, hence
\[
    u_{i,n}(F)=\frac{\lvert f_i(F) \rvert}{\lVert f_i(F) \rVert},
    \qquad
    u_{j,n}(F)=\frac{\lvert f_j(F) \rvert}{\lVert f_j(F) \rVert},
\]
so the penalty terms vanish.
Therefore, for all such $n$,
\[
\begin{aligned}
    \psi_{i,j,n}(F)
    &= \lVert u_{i,n}(F) + \varepsilon\,u_{j,n}(F) \rVert - 1 \\
    &\le \lVert x + \varepsilon y \rVert - 1
       + \lVert u_{i,n}(F) - x \rVert
       + \varepsilon\,\lVert u_{j,n}(F) - y \rVert \\
    &< \sigma_F(\varepsilon) + (2+\varepsilon)\eta.
\end{aligned}
\]
Hence
$\limsup_n \sigma^{(n)}_F(\varepsilon) \le \sigma_F(\varepsilon)$.
Combined with the opposite inequality, this proves
$\sigma^{(n)}_F(\varepsilon)\downarrow \sigma_F(\varepsilon)$.

Therefore $F \mapsto \sigma_F(\varepsilon)$ is upper semicontinuous on
$\Subl(E)\setminus\{\{0\}\}$.
Since changing the value of a Baire class~$1$ function at a single point preserves
Baire class~$1$, and $\sigma_{\{0\}}(\varepsilon)=0$ by convention,
$F \mapsto \sigma_F(\varepsilon)$ is Baire class~$1$ on all of $\Subl(E)$.
Consequently,
\[
    \{F \colon \sigma_F(\varepsilon) > 0\}
    = \bigcup_{m \ge 1} \{F \colon \sigma_F(\varepsilon) \ge 1/m\}
\]
is $F_\sigma$.

\smallskip
\emph{Lower semicontinuous approximants for $\rho_F(\tau)$.}
Fix a rational $\tau \in (0, 1]$.  An analogous reduction, valid for every
$F \ne \{0\}$, gives
\begin{equation}\label{eq:rho-unit-reduction}
    \rho_F(\tau)
    = \sup\bigl\{\lVert x \vee \tau\, y \rVert - 1
      \colon x, y \in S(F_+)\bigr\}.
\end{equation}
For $i, j \in \omega$ and $n \ge 1$ define
\[
    \chi_{i,j,n}(F)
    := \max\Bigl\{0,\;
    \lVert u_{i,n}(F) \vee \tau\, u_{j,n}(F) \rVert - 1\Bigr\}.
\]
Each $\chi_{i,j,n}$ is continuous (since $\vee$ is continuous by
Lemma~\ref{lem:lattice-ops-lipschitz}\,(ii), the norm is continuous, and the
maximum with~$0$ is continuous).  Put
\[
    \rho^{(n)}_F(\tau) := \sup_{i,j \in \omega} \chi_{i,j,n}(F).
\]
Being the pointwise supremum of countably many \emph{continuous} functions,
$\rho^{(n)}_\bullet(\tau)$ is lower semicontinuous, hence Baire class~$1$.
Moreover, $\rho^{(n)}_F(\tau)$ is non-decreasing in~$n$: since the scalar
multipliers defining $u_{j,n}(F)$ are non-decreasing in~$n$, so are
$u_{j,n}(F)$ in the lattice order, and the lattice join and norm are
order-preserving on $E^+$.

We now show that
$\rho^{(n)}_F(\tau) \uparrow \rho_F(\tau)$ for every $F \ne \{0\}$.

For fixed $F$ and $j$, the vectors $u_{j,n}(F)$ increase in the lattice order with $n$,
because the scalar multipliers defining them are non-decreasing.
Hence, for fixed $i,j$,
\[
    u_{i,n}(F) \vee \tau\,u_{j,n}(F)
    \preceq
    u_{i,n+1}(F) \vee \tau\,u_{j,n+1}(F),
\]
and therefore
$\chi_{i,j,n}(F) \le \chi_{i,j,n+1}(F)$.
Taking suprema gives that $\rho^{(n)}_F(\tau)$ is non-decreasing in $n$.

Also, $\rho^{(n)}_F(\tau) \le \rho_F(\tau)$.
Indeed, if $a,b \in F_+$ satisfy $\lVert a \rVert,\lVert b \rVert \le 1$ and
both are non-zero, write
$a=\alpha x$, $b=\beta y$ with $\alpha,\beta \in (0,1]$ and
$x,y \in S(F_+)$.
Then $a \preceq x$ and $\tau b \preceq \tau y$, so
\[
    a \vee \tau b \preceq x \vee \tau y.
\]
By norm monotonicity,
\[
    \lVert a \vee \tau b \rVert - 1
    \le \lVert x \vee \tau y \rVert - 1
    \le \rho_F(\tau).
\]
If one of $a,b$ is zero, then
$\chi_{i,j,n}(F)=0 \le \rho_F(\tau)$.
Applying this to
$a=u_{i,n}(F)$ and $b=u_{j,n}(F)$ yields
$\rho^{(n)}_F(\tau) \le \rho_F(\tau)$.

Conversely, let $\eta > 0$.
By~\eqref{eq:rho-unit-reduction}, choose $x,y \in S(F_+)$ such that
\[
    \lVert x \vee \tau y \rVert - 1 > \rho_F(\tau) - \eta.
\]
As above, choose indices $i,j$ such that
\[
    \left\lVert
    \frac{\lvert f_i(F) \rvert}{\lVert f_i(F) \rVert} - x
    \right\rVert < \eta,
    \qquad
    \left\lVert
    \frac{\lvert f_j(F) \rvert}{\lVert f_j(F) \rVert} - y
    \right\rVert < \eta.
\]
For all sufficiently large $n$, the truncation is inactive at $f_i(F)$ and $f_j(F)$, so
\[
    u_{i,n}(F)=\frac{\lvert f_i(F) \rvert}{\lVert f_i(F) \rVert},
    \qquad
    u_{j,n}(F)=\frac{\lvert f_j(F) \rvert}{\lVert f_j(F) \rVert}.
\]
By continuity of $\vee$ and of the norm,
\[
    \chi_{i,j,n}(F) > \rho_F(\tau) - 2\eta
\]
for all sufficiently large $n$.
Hence
$\liminf_n \rho^{(n)}_F(\tau) \ge \rho_F(\tau)$.
Combined with the opposite inequality, this proves
$\rho^{(n)}_F(\tau)\uparrow \rho_F(\tau)$.

Thus $F \mapsto \rho_F(\tau)$ is lower semicontinuous on
$\Subl(E)\setminus\{\{0\}\}$.
Since $\rho_{\{0\}}(\tau)=0$ by convention and changing the value at a single
point preserves Baire class~$1$, $F \mapsto \rho_F(\tau)$ is Baire class~$1$
on all of $\Subl(E)$.
Therefore, for every rational $r$, the set
$\{F \colon \rho_F(\tau) < r\}$ is $F_\sigma$.

\medskip
\emph{Proof of~\ref{kurc:UM}.}
Kurc shows that $\sigma_X(\varepsilon)$ (and hence $\delta_X(\varepsilon)$) is
non-decreasing in $\varepsilon$; see~\cite{Kurc1993Dual}.  Therefore
$\delta_F(\varepsilon) > 0$ for all $\varepsilon \in (0, 1]$ if and only if
$\sigma_F(1/n) > 0$ for all $n \ge 1$.  Hence
\[
    \mathrm{UM}(E) = \bigcap_{n \ge 1}
    \{F \in \Subl(E) \colon \sigma_F(1/n) > 0\}.
\]
By~\ref{kurc:baire1}, each set $\{F \colon \sigma_F(1/n) > 0\}$ is $F_\sigma$
($\boldsymbol{\Sigma}^0_2$).  A countable intersection of
$\boldsymbol{\Sigma}^0_2$ sets is $\boldsymbol{\Pi}^0_3$, so
$\mathrm{UM}(E) \in \boldsymbol{\Pi}^0_3(\Subl(E))$.

\medskip
\emph{Proof of~\ref{kurc:oUSm}.}
Kurc observes that for fixed $x, y \succeq \mathbf{0}_X$ with unit norm, the
function $t \mapsto \lVert x \vee t y \rVert$ is convex and satisfies
$\lVert x \vee 0 \cdot y \rVert = \lVert x \rVert = 1$.  Setting
$g(t) = \lVert x \vee t y \rVert - 1$, we have $g(0) = 0$ and $g$ convex,
so for $0 < s < t$ convexity gives
$g(s) \le (s/t)\, g(t) + (1 - s/t)\, g(0) = (s/t)\, g(t)$, whence
$g(s)/s \le g(t)/t$.  Thus $t \mapsto (\lVert x \vee t y \rVert - 1) / t$ is
non-decreasing on $(0, 1]$.  Taking the supremum over all such pairs preserves
this monotonicity, so $\tau \mapsto \rho_F(\tau) / \tau$ is non-decreasing
on $(0, 1]$.

Therefore $\rho_F(\tau) / \tau \to 0$ as $\tau \downarrow 0$ if and only if
for every $k \ge 1$ there exists $n \ge 1$ such that
$\rho_F(1/n) / (1/n) < 1/k$, i.e.\ $\rho_F(1/n) < 1/(kn)$.  Hence
\[
    \mathrm{(o)\text{-}USm}(E)
    = \bigcap_{k \ge 1} \bigcup_{n \ge 1}
      \bigl\{F \in \Subl(E) \colon \rho_F(1/n) < 1/(kn)\bigr\}.
\]
By~\ref{kurc:baire1}, each set $\{F \colon \rho_F(1/n) < 1/(kn)\}$ is
$F_\sigma$.  A countable union of $F_\sigma$ sets is $F_\sigma$, so
$\bigcup_{n \ge 1} \{F \colon \rho_F(1/n) < 1/(kn)\}$ is $F_\sigma$
($\boldsymbol{\Sigma}^0_2$).  A countable intersection of
$\boldsymbol{\Sigma}^0_2$ sets is $\boldsymbol{\Pi}^0_3$, so
$\mathrm{(o)\text{-}USm}(E) \in \boldsymbol{\Pi}^0_3(\Subl(E))$.

\medskip
\emph{Proof of~\ref{kurc:STM}.}
It suffices to show that the complement of $\mathrm{STM}(E)$ is analytic
($\boldsymbol{\Sigma}^1_1$).  Define
\begin{multline*}
    \mathcal{A} = \bigl\{(F, x, y) \in \Subl(E) \times E \times E \colon
    x \in F,\; y \in F,\; x \in E^+,\; y \in E^+,\\
    y \ne 0,\; y \preceq x,\;
    \lVert x - y \rVert = \lVert x \rVert\bigr\}.
\end{multline*}
We verify that $\mathcal{A}$ is Borel in
$\Subl(E) \times E \times E$.  The conditions $x \in E^+$ and $y \in E^+$ each
define closed subsets of~$E$ (since $E^+$ is norm-closed).  The condition
$y \preceq x$ (equivalently $x - y \in E^+$) is closed.  The condition
$\lVert x - y \rVert = \lVert x \rVert$ is closed (the norm is continuous).
The condition $y \ne 0$ is open.  The membership conditions $x \in F$ and
$y \in F$ are each closed in $E \times \Subl(E)$ by~\ref{A3}.  Hence
$\mathcal{A}$ is a Borel subset of the product.

The complement of $\mathrm{STM}(E)$ is the projection
$\pi_1(\mathcal{A}) \subseteq \Subl(E)$, where
$\pi_1(F, x, y) = F$.  The continuous image (projection) of a Borel subset of a
product of Polish spaces is analytic.  Therefore
$\Subl(E) \setminus \mathrm{STM}(E)$ is $\boldsymbol{\Sigma}^1_1$, and so
$\mathrm{STM}(E)$ is $\boldsymbol{\Pi}^1_1$.
\end{proof}

\begin{remark}\label{rem:kurc-topology-transfer}
The Borel/analytic/coanalytic status of the sets in
Theorem~\ref{thm:kurc-borel} does not depend on the choice of admissible
topology, since any two admissible topologies generate the same Effros--Borel
$\sigma$-algebra.  However, the Borel \emph{rank} may shift: since the
identity map between two admissible topologies is Baire class~$1$,
a $\boldsymbol{\Pi}^0_3$ set in the Wijsman topology remains Borel in any
other admissible topology, and is at most $\boldsymbol{\Pi}^0_4$ there.
\end{remark}

\begin{remark}\label{rem:kurc-dual}
Kurc's (o)-Sm property is defined using the subdifferential
$\partial_+\lVert x \rVert \subseteq X^*_+$, which involves the dual
lattice.  Without a Borel coding of the dual unit sphere as a functor on
$\Subl(E)$, we do not presently determine the Borel rank of (o)-Sm.  However,
Kurc's duality results~\cite{Kurc1993Dual} imply that once a functorial
Borel dual-coding is established, the complexity of STM and (o)-Sm would
transfer between primal and dual codes.
\end{remark}

\subsection{Strict monotonicity: a test family and the completeness problem}\label{subsec:STM-complete}

The preceding theorem gives the coanalytic upper bound for strict monotonicity.
A natural next question is whether this upper bound is sharp.  The following
simple observation is a useful source of non-strictly monotone examples, but it
also shows why a direct reduction through compact hyperspaces is too naive.

\begin{lemma}\label{lem:C(K)-STM}
Let $K$ be a non-empty compact Hausdorff space.  Then the Banach lattice $C(K)$
(with the supremum norm and pointwise order) is strictly monotone if and only if
$K$ is a singleton.
\end{lemma}

\begin{proof}
If $K$ is a singleton, then $C(K)\cong\R$, which is strictly monotone.

Conversely, assume that $K$ contains two distinct points $p\ne q$.  By
Urysohn's lemma there exists $y\in C(K)$ with $0\le y\le 1$, $y(p)=0$, and
$y(q)=1$.  Let $x\equiv 1$.  Then $0\prec y\preceq x$, but $x-y$ attains the
value $1$ at~$p$, so
\[
    \|x-y\|_\infty=1=\|x\|_\infty .
\]
Thus $C(K)$ is not strictly monotone.
\end{proof}

\begin{remark}\label{rem:STM-complete-caution}
Lemma~\ref{lem:C(K)-STM} suggests trying to reduce well-foundedness of trees to
strict monotonicity by assigning to a tree $T$ a compact space $K_T$ which is a
singleton exactly when $T$ is well founded.  Such an argument cannot proceed via
a Borel map $T\mapsto K_T$ into the hyperspace $\mathcal F(\Delta)$: the set of
singleton compacta is Borel in $\mathcal F(\Delta)$, whereas the classical
well-founded-tree set is non-Borel.  One must also use trees on
$\N^{<\omega}$, not on $2^{<\omega}$, for the classical
$\boldsymbol{\Pi}^1_1$-complete well-foundedness problem; on finitely branching
binary trees, well-foundedness is Borel by K\"onig's lemma.

Consequently, a proof that $\mathrm{STM}(\mathcal C)$ is
$\boldsymbol{\Pi}^1_1$-hard would need a genuinely Banach-lattice-theoretic
coding, rather than only the elementary compact-space test above.
\end{remark}

\begin{problem}\label{prob:STM-complete}
Is $\mathrm{STM}(\mathcal C)\subseteq\Subl(\mathcal C)$
$\boldsymbol{\Pi}^1_1$-complete?  More generally, for which separable ambient
Banach lattices $E$ is $\mathrm{STM}(E)$ coanalytic-complete in $\Subl(E)$?
\end{problem}

\subsection{Open questions}

Determining the Borel complexity of the classes of lattices possessing the Fatou,
Nakano, strong Nakano, or projectivity properties within our coding spaces, as
well as characterising metrisable SICK compacta, are natural open problems.
Whether the $\boldsymbol{\Pi}^0_3$ bounds for UM and (o)-USm in
Theorem~\ref{thm:kurc-borel} are sharp (i.e.\ $\boldsymbol{\Pi}^0_3$-complete)
and the determination of the exact complexity of (o)-Sm are also of interest.
Problem~\ref{prob:STM-complete} asks whether the coanalytic upper bound for
strict monotonicity is sharp in the universal coding space.

\begin{remark}[Isometric lattice isomorphism]\label{rem:isom-relation}
The relation of \emph{isometric lattice isomorphism} on $\Subl(\mathcal{C})$,
defined by $Y_1 \sim Y_2$ if and only if $Y_1$ and $Y_2$ are isometrically
lattice-isomorphic, is analytic ($\boldsymbol{\Sigma}^1_1$).

Indeed, rather than quantifying over the non-separable operator space
$\mathcal{L}(\mathcal{C})$, one may code an operator by its values on a fixed
dense sequence in $\mathcal{C}$; this yields a Polish parameter space
$\mathcal{C}^\omega$ of operator codes.  Using the continuous selectors from
Theorem~\ref{thm:selection}, one can express in a Borel way that such a code
determines a surjective linear isometry $T \colon Y_1 \to Y_2$ satisfying
$T(\lvert x \rvert) = \lvert T x \rvert$ on a dense set, hence on all of $Y_1$
by continuity of the lattice operations.  Therefore $\sim$ is the projection of
a Borel subset of
$\Subl(\mathcal{C}) \times \Subl(\mathcal{C}) \times \mathcal{C}^\omega$.

Determining whether $\sim$ is Borel remains a natural structural problem for
this coding space.
\end{remark}

\begin{remark}[Complexity of metrisable SICK compacta]\label{rem:SICK-complexity}
Let $\mathbb H=[0,1]^\omega$ be the Hilbert cube and equip
$\mathcal{F}(\mathbb H)$ with the Vietoris topology.  Since every compact
metrisable space is homeomorphic to a closed subset of $\mathbb H$, this is the
natural hyperspace for the metrisable SICK question.  Define
\[
    \mathrm{SICK}_{\mathrm{met}}
    := \bigl\{K \in \mathcal{F}(\mathbb H) \colon K \text{ is a metrisable SICK
    compactum}\bigr\}.
\]
Using $\mathcal{F}(\Delta)$ instead would restrict the problem further to
zero-dimensional compacta, because closed subsets of the Cantor set are
zero-dimensional.  A treatment of all, possibly non-metrisable, SICK compacta
would require a different standard Borel coding of compact Hausdorff spaces, or
an indirect coding through the separable Banach lattices and principal ideals
that produce their structure spaces.

The characterisation in~\cite{AMRT2022Sick} suggests that
$\mathrm{SICK}_{\mathrm{met}}$ should be at most analytic, but the exact level in
the projective hierarchy remains open.
\end{remark}

% The resolution of
% Problem~\ref{prob:coding-equiv} is also of particular interest; it would
% establish a definitive link between the sublattice and ideal coding spaces.

\subsection*{Acknowledgements}
The author gratefully acknowledges support received from
NCN Sonata-Bis~13 (2023/50/E/ST1/00067).


%=============================================================================
% Bibliography
%=============================================================================

\begin{thebibliography}{99}

\bibitem{aviles2018free}
A.~Avil\'es, J.~Rodr\'\i guez, and P.~Tradacete,
\textit{The free Banach lattice generated by a Banach space},
J.~Funct.~Anal.\ \textbf{274} (2018), no.~10, 2955--2977.

\bibitem{AMRT2022Sick}
A.~Avil\'es, G.~Mart\'\i nez-Cervantes, A.~Rueda Zoca, and P.~Tradacete,
\textit{Compact spaces associated to separable Banach lattices},
Banach J.~Math.~Anal.\ \textbf{16} (2022), no.~4, Art.~67.

\bibitem{Boss2002Cod}
B.~Bossard,
\textit{A coding of separable Banach spaces.\ Analytic and coanalytic families
of Banach spaces},
Fund.~Math.\ \textbf{172} (2002), no.~2, 117--152.

\bibitem{de2015free}
B.~de~Pagter and A.~W.~Wickstead,
\textit{Free and projective Banach lattices},
Proc.~Roy.~Soc.~Edinburgh Sect.~A \textbf{145} (2015), no.~1, 105--143.

\bibitem{Dod2010Ban}
G.~Dodos,
\textit{Banach Spaces and Descriptive Set Theory: Selected Topics},
Lecture Notes in Mathematics, vol.~1993, Springer-Verlag, 2010.

\bibitem{Frem1974Ban}
D.~H.~Fremlin,
\textit{Tensor products of Banach lattices},
Math.~Ann.\ \textbf{211} (1974), 87--106.

\bibitem{God2018Cod}
G.~Godefroy and J.~Saint-Raymond,
\textit{Descriptive complexity of some isomorphism classes among Banach spaces},
J.~Funct.~Anal.\ \textbf{275} (2018), no.~4, 1008--1022.

\bibitem{Kan2026Pol}
T.~Kania,
\textit{Polish spaces for countable and separable structures through quotient
encodings},
arXiv:2604.15843 [math.LO], 2026,
\href{https://doi.org/10.48550/arXiv.2604.15843}{doi:10.48550/arXiv.2604.15843}.

\bibitem{Kurc1993Dual}
W.~Kurc,
\textit{A dual property to uniform monotonicity in Banach lattices},
Collect.\ Math.\ \textbf{44} (1993), 155--165.

\bibitem{LLOT2019Uni}
D.~H.~Leung, L.~Li, T.~Oikhberg, and M.~A.~Tursi,
\textit{Separable universal Banach lattices},
Israel J.~Math.\ \textbf{230} (2019), 141--152.

\bibitem{Nie1982Ideal}
N.~J.~Nielsen,
\textit{The ideal property of tensor products of Banach lattices with
applications to the local structure of spaces of absolutely summing operators},
Studia Math.\ \textbf{74} (1982), no.~3, 247--272.

\bibitem{Pug2007Fremlin}
D.~Puglisi,
\textit{Fremlin tensor products of Banach lattices},
Quaestiones Math.\ \textbf{30} (2007), no.~1, 45--56.

\bibitem{Scha1974Ban}
H.~H.~Schaefer,
\textit{Banach Lattices and Positive Operators},
Grundlehren der mathematischen Wissenschaften, vol.~215, Springer-Verlag, 1974.

\bibitem{Wick2007nak}
A.~W.~Wickstead,
\textit{Banach lattices with the Nakano property},
Positivity \textbf{11} (2007), no.~2, 329--339.

\end{thebibliography}

\end{document}
